% T12 C/D, chapitre 1 : Arithmétique. Source : le livre « Terminales scientifiques » (Arithmetique/Lesona, Arithmetique/Exo, Arithmetique/Valiny).
% Les exercices tirés d’autres ouvrages ne sont pas repris (décision du 2 novembre 2026).

\TCDLessonStyleOn
\providecommand{\congru}[3]{#1 \equiv #2 \quad \mathrm{mod}\ #3}

\chapter{Arithm\'etique}
\begin{defi}[Divisibilit\'e] 
\index{Divisibilite@Divisibilit\'e}
Soient $a$ et $b$ deux entiers naturels.
On dit que \emph{$a$ divise $b$} ou que \emph{$a$ est un diviseur de $b$} s'il existe un entier naturel $q$ tel que $b = aq$.
On dit aussi que \emph{$b$ est un multiple de $a$} ou encore que \emph{$b$ est divisible par $a$}.

L'ensemble des diviseurs de $a$ est not\'e $D_{(a)}$ ou $D_a$.
\index{Divisibilite@Divisibilit\'e!Diviseur}
\end{defi}
            %
      \section{Nombres premiers}
\index{Nombres premiers}
            \begin{defi}[Nombre premier]
\index{Premiers entre eux}
                  \begin{enumerate}
                        \item Un nombre $n\in\mathbb{N}^*$ est dit \emph{premier} 
                              s'il admet exactement deux diviseurs positifs.
                        \item Deux nombres entiers relatifs $a$ et $b$ sont dits \emph{premiers entre eux} s'ils
                              n'ont pas de diviseur commun positif autre que $1$.
                              %si $\pgcd{a}{b} = 1$.
                  \end{enumerate}
            \smallnota{Le nombre $1$ \textbf{n'}est \textbf{pas} premier,
                        le plus petit nombre premier est $2$.}
            \end{defi}
            %
            \begin{prop}
                  \begin{itemize}
                        \item Tout nombre entier naturel $n > 1$ qui n'est pas premier 
                              admet au moins un diviseur premier $p \leq \sqrt{n}$.
                        \item Tout nombre entier naturel $n \geq 2$ peut se d\'ecomposer en produit de facteurs premiers,
                              cette d\'ecomposition est unique \`a l'ordre pr\`es.
                  \end{itemize}
            %
                  \begin{ex}
                        D\'ecomposer $90$ et $140$ en produits de facteurs premiers.
                        \[
                                    \begin{array}{r|l}
                                           90 & 2 \\
                                           45 & 3 \\
                                           15 & 3 \\
                                            5 & 5 \\
                                            1 &   \\
                                    \end{array}
                                    \qquad
                                    \begin{array}{r|l}
                                          140 & 2 \\
                                           70 & 2 \\
                                           35 & 5 \\
                                            7 & 7 \\
                                            1 &   \\
                                    \end{array}
                        \]
                                    Ainsi, $90 = 2 \cdot 3^2 \cdot 5$ et $140 = 2^2 \cdot 5 \cdot 7$.
                  \end{ex}
            %
            Etant donn\'e un nombre entier naturel $n > 2$,
            pour savoir si $n$ est premier,
            il n'est pas n\'ecessaire de v\'erifier qu'aucun nombre $q < n$ ne divise $n$.
            Il suffit de faire la v\'erification pour les nombres premiers inf\'erieurs ou \'egaux \`a $\sqrt{n}$.
            %
                  \begin{ex}
                        V\'erifions si $139$ est premier.
                        
                        Commen\c{c}ons par trouver $\sqrt{139} \simeq 11,8$.
                        Listons tous les nombres premiers inf\'erieurs \`a $11,8$ : 
                              2, 3, 5, 7 et 11.
                        En effectuant la division de $139$ par ces nombres, on constate qu'aucun d'eux ne divise $139$.
                        On conclut que $139$ est premier.
                  \end{ex}
            %
                  \begin{exe}\label{exe:arith-1}
                        Lesquels des nombres suivants sont premiers : 
                              527, 769, 1249 et 
                              57.
%                              \footnote{
%                                          Le cel\`ebre math\'ematicien Alexander Grothendieck,
%                                          lors d'une conf\'erence,
%                                          exposa de mani\`ere un peut trop th\'eorique comme \`a son habitude.
%                                          Un membre de l'audience demanda alors un exemple concret,
%                                          c'est-\`a-dire utilisant un nombre.
%                                          Grothendieck proposa alors de prendre ``un nombre premier,
%                                          disons 57 \ldots ''.
%                                          }
                  \end{exe}
            \end{prop}
            %
            \begin{thm}
                  Il y a une infinit\'e de nombres premiers.
            \end{thm}
            %
            \subsection*{Crible d'Eratosthenes}
\index{Eratosthenes!Crible d'}
                  \begin{figure}[h!]
                  \centering
                        \begin{tikzpicture}[scale=.8]
\clip[draw] (0,-2) rectangle (10,10);
%
\begin{scope}[opacity=.25]
%
\foreach \x in {2,4,...,10}
      \filldraw[color = loko1] (\x-1,10) rectangle (\x,9);
\foreach \x in {12,14,...,20}
      \filldraw[color = loko1] (\x-11,9) rectangle (\x-10,8);
\foreach \x in {22,24,...,30}
      \filldraw[color = loko1] (\x-21,8) rectangle (\x-20,7);
\foreach \x in {32,34,...,40}
      \filldraw[color = loko1] (\x-31,7) rectangle (\x-30,6);
\foreach \x in {42,44,...,50}
      \filldraw[color = loko1] (\x-41,6) rectangle (\x-40,5);
\foreach \x in {52,54,...,60}
      \filldraw[color = loko1] (\x-51,5) rectangle (\x-50,4);
\foreach \x in {62,64,...,70}
      \filldraw[color = loko1] (\x-61,4) rectangle (\x-60,3);
\foreach \x in {72,74,...,80}
      \filldraw[color = loko1] (\x-71,3) rectangle (\x-70,2);
\foreach \x in {82,84,...,90}
      \filldraw[color = loko1] (\x-81,2) rectangle (\x-80,1);
\foreach \x in {92,94,...,100}
      \filldraw[color = loko1] (\x-91,1) rectangle (\x-90,0);
\foreach \x in {102,104,...,110}
      \filldraw[color = loko1] (\x-101,0) rectangle (\x-100,-1);
\foreach \x in {112,114,...,120}
      \filldraw[color = loko1] (\x-111,-1) rectangle (\x-110,-2);
%
\foreach \x in {3,6,...,10}
      \filldraw[color = loko3] (\x-1,10) rectangle (\x,9);
\foreach \x in {12,15,...,20}
      \filldraw[color = loko3] (\x-11,9) rectangle (\x-10,8);
\foreach \x in {21,24,...,30}
      \filldraw[color = loko3] (\x-21,8) rectangle (\x-20,7);
\foreach \x in {30,33,...,40}
      \filldraw[color = loko3] (\x-31,7) rectangle (\x-30,6);
\foreach \x in {42,45,...,50}
      \filldraw[color = loko3] (\x-41,6) rectangle (\x-40,5);
\foreach \x in {51,54,...,60}
      \filldraw[color = loko3] (\x-51,5) rectangle (\x-50,4);
\foreach \x in {60,63,...,70}
      \filldraw[color = loko3] (\x-61,4) rectangle (\x-60,3);
\foreach \x in {72,75,...,80}
      \filldraw[color = loko3] (\x-71,3) rectangle (\x-70,2);
\foreach \x in {81,84,...,90}
      \filldraw[color = loko3] (\x-81,2) rectangle (\x-80,1);
\foreach \x in {90,93,...,100}
      \filldraw[color = loko3] (\x-91,1) rectangle (\x-90,0);
\foreach \x in {102,105,...,110}
      \filldraw[color = loko3] (\x-101,0) rectangle (\x-100,-1);
\foreach \x in {111,114,...,120}
      \filldraw[color = loko3] (\x-111,-1) rectangle (\x-110,-2);
%
\foreach \x in {5,10}
      \filldraw[color = loko2] (\x,10) rectangle (\x-1,9);
\foreach \x in {15,20}
      \filldraw[color = loko2] (\x-10,9) rectangle (\x-11,8);
\foreach \x in {25,30}
      \filldraw[color = loko2] (\x-20,8) rectangle (\x-21,7);
\foreach \x in {35,40}
      \filldraw[color = loko2] (\x-30,7) rectangle (\x-31,6);
\foreach \x in {45,50}
      \filldraw[color = loko2] (\x-40,6) rectangle (\x-41,5);
\foreach \x in {55,60}
      \filldraw[color = loko2] (\x-50,5) rectangle (\x-51,4);
\foreach \x in {65,70}
      \filldraw[color = loko2] (\x-60,4) rectangle (\x-61,3);
\foreach \x in {75,80}
      \filldraw[color = loko2] (\x-70,3) rectangle (\x-71,2);
\foreach \x in {85,90}
      \filldraw[color = loko2] (\x-80,2) rectangle (\x-81,1);
\foreach \x in {95,100}
      \filldraw[color = loko2] (\x-90,1) rectangle (\x-91,0);
\foreach \x in {105,110}
      \filldraw[color = loko2] (\x-100,0) rectangle (\x-101,-1);
\foreach \x in {115,120}
      \filldraw[color = loko2] (\x-110,-1) rectangle (\x-111,-2);
%
\foreach \x in {14}
      \filldraw[color = loko5] (\x-10,9) rectangle (\x-11,8);
\foreach \x in {21,28}
      \filldraw[color = loko5] (\x-20,8) rectangle (\x-21,7);
\foreach \x in {35}
      \filldraw[color = loko5] (\x-30,7) rectangle (\x-31,6);
\foreach \x in {42,49}
      \filldraw[color = loko5] (\x-40,6) rectangle (\x-41,5);
\foreach \x in {56}
      \filldraw[color = loko5] (\x-50,5) rectangle (\x-51,4);
\foreach \x in {63,70}
      \filldraw[color = loko5] (\x-60,4) rectangle (\x-61,3);
\foreach \x in {77}
      \filldraw[color = loko5] (\x-70,3) rectangle (\x-71,2);
\foreach \x in {84}
      \filldraw[color = loko5] (\x-80,2) rectangle (\x-81,1);
\foreach \x in {91,98}
      \filldraw[color = loko5] (\x-90,1) rectangle (\x-91,0);
\foreach \x in {105}
      \filldraw[color = loko5] (\x-100,0) rectangle (\x-101,-1);
\foreach \x in {112,119}
      \filldraw[color = loko5] (\x-110,-1) rectangle (\x-111,-2);
%
\end{scope}
%
\textbf{
%
\foreach \x in {1,2,...,10}
      \draw[] (\x-.5,9.5) node{\x};
\foreach \x in {11,12,...,20}
      \draw (\x-10.5,8.5) node{\x};
\foreach \x in {21,22,...,30}
      \draw (\x-20.5,7.5) node{\x};
\foreach \x in {31,32,...,40}
      \draw (\x-30.5,6.5) node{\x};
\foreach \x in {41,42,...,50}
      \draw (\x-40.5,5.5) node{\x};
\foreach \x in {51,52,...,60}
      \draw (\x-50.5,4.5) node{\x};
\foreach \x in {61,62,...,70}
      \draw (\x-60.5,3.5) node{\x};
\foreach \x in {71,72,...,80}
      \draw (\x-70.5,2.5) node{\x};
\foreach \x in {81,82,...,90}
      \draw (\x-80.5,1.5) node{\x};
\foreach \x in {91,92,...,100}
      \draw (\x-90.5,.5) node{\x};
\foreach \x in {101,102,...,110}
      \draw (\x-100.5,-.5) node{\x};
\foreach \x in {111,112,...,120}
      \draw (\x-110.5,-1.5) node{\x};
%
}
%
\filldraw[opacity=.25] (0,10) rectangle (1,9);
%\filldraw[] (.5,9.5) node[] {1};
\filldraw[color = loko1] (1,10) rectangle (2,9);
\filldraw[] (1.5,9.5) node[fill=white] {2};
\filldraw[color = loko3] (2,10) rectangle (3,9);
\filldraw[] (2.5,9.5) node[fill=white] {3};
\filldraw[color = loko2] (4,10) rectangle (5,9);
\filldraw[] (4.5,9.5) node[fill=white] {5};
\filldraw[color = loko5] (6,10) rectangle (7,9);
\filldraw[] (6.5,9.5) node[fill=white] {7};
%
\draw[line width=1] (0,-2) grid (10,10);
\end{tikzpicture}
                  \caption{Crible d'Eratosthenes pour les nombres de $1$ \`a $120$.
%                           Les nombres premiers entre $1$ et $120$ sont : 
%                           2, 3, 5, 7, 11, 13, 17, 19, 23, 29, 31, 37, 41, 43, 47, 53, 59, 
%                           61, 67, 71, 73, 79, 83, 89, 97, 101, 103, 107, 109 et 113.
                              }
                  \label{fig: Eratosthenes}
                  \end{figure}
            Le crible d'Eratosthenes est une m\'ethode pour trouver tous les nombres premiers jusqu'\`a un rang $n$ donn\'e.
            On commence par \'ecrire tous les nombres entiers de $1$ \`a $n$ dans une table 
            contenant $\lfloor \sqrt{n} \rfloor$%
            \footnote{ $\lfloor X \rfloor$ ou $\mathrm{E}\left( X \right)$ signifie 
                       \emph{partie enti\`ere de $X$} ;
\index{Partie entiere@Partie enti\`ere}
                       c'est-\`a-dire le plus grand entier inf\'erieur ou \'egal \`a $X$.
%                       Par exemple :
%                       \begin{align*}
%                             \lfloor 4,42 \rfloor  & = 4 &
%                             \lfloor \pi \rfloor  & = 3 \\
%                             \lfloor - 4,42 \rfloor  & = -5 \text{ (car $-5 \leq -4,42 < -4$) }   &
%                             \lfloor - \pi \rfloor & = -4 \text{ (car $-4 \leq -\pi < -3$) }      \;.
%                       \end{align*}
                     }
            colonnes.
            On raye le nombre $1$ puisqu'il n'est pas premier.
            Le prochain nombre non-ray\'es est premier, on le prend et on raye tous ses multiples.
            On recommence le proc\'ed\'e sur toute la premi\`ere ligne.
            On s'arr\^ete apr\`es avoir trait\'e les nombres premiers inf\'erieurs ou \'egaux \`a $\sqrt{n}$.
            \`A la fin de l'algorithme, tous les nombres non ray\'es sont premiers.
%                  
\section{Diviseurs et multiples d'un nombre entier relatif}
\begin{defi}[Divisibilit\'e] 
\index{Divisibilite@Divisibilit\'e!divise}
Soient $a$ et $b$ deux nombres entiers relatifs.
On dit que \emph{$a$ divise $b$}, et on note \defbox{$a | b$}, s'il existe un nombre entier relatif $q$ tel que $b = q a$.
On dit aussi que \emph{$b$ est un multiple de $a$}; \index{multiple}
le nombre $q$ est appel\'e \emph{quotient} de $b$ par $a$.
\end{defi}
%
\begin{prop}
\label{prop:divisibilite}
Soient $a$, $b$ et $c$ trois nombres entiers relatifs non-nuls.
      \begin{enumerate}
      \item Si $a|b$ et $b|a$ alors $a = b$ ou $a = -b$.
      \item \begin{enumerate}
            \item $a|a$,
            \item $(a|b \text{ et } b|c) \Rightarrow a|c$,
            \item $(a|b \text{ et } b|a) \Rightarrow |a| = |b|$.
            \end{enumerate}
      \item \begin{enumerate}
            \item Si $a|b$ alors $a|bc$,
            \item si $a|b$ alors $ac|bc$,
            \item si $a|b$ et $a|c$ alors $a|(\beta b + \gamma c)$, pour tout $(\beta, \gamma) \in \mathbb{Z}^2$.
            \end{enumerate}
      \item $1|a$ et $a|0$.
      \item $b|a \text{ et } a\neq 0 \Rightarrow 0 < |b| \leq |a|$.
      \end{enumerate}
%
%\begin{proof}
%Voir preuve~\ref{prop:divisibilite}
%\end{proof}
\end{prop}

\paragraph{Division euclidienne dans $\mathbb{Z}$} 
\index{Division euclidienne}
\begin{thm}[Admis]
\label{thm:division euclidienne}
Soit $(a,b) \in \mathbb{Z} \times \mathbb{Z}^*$.
\par\smallskip
Il existe alors un unique couple d'entiers relatifs $(q,r)$ tel que
\par\noindent
      \thmbox{
      $ a = b q + r \text{ et } 0 \leq r < |b|      \;. $
      }
\begin{description}
      \item[$q$] est le quotient de la division de $a$ par $b$,      \index{Division euclidienne!quotient}
      \item[$r$] est le reste de la division de $a$ par $b$.      \index{Division euclidienne!reste}
\end{description}
$b|a$ si et seulement si $r=0$.
      %
      \begin{ex}
      Prenons $a = -17$ et $b = -4$; on a $q = 5$ et $r = 3$.
            On v\'erifie que $-17 = (-4) \cdot 5 + 3$.
      \end{ex}
\end{thm}
%
\section{Num\'eration \`a base $b$}
\index{Numeration@Num\'eration}
      Rappelons que dans l'\'ecriture usuelle d'un nombre entier, on utilise la base $10$.
      On se sert des symboles : $0$, $1$, $2$, $3$, $4$, $5$, $6$, $7$, $8$ et $9$ (il y a dix symboles).
      
      Soit $N$ un entier naturel et
            \defbox{$ a_n a_{n-1} \dots a_2 a_1 a_0 $}
      o\`u $a_i \in \{ 0, 1, 2, \dots , 9 \}$ est son \'ecriture d\'ecimale (c'est-\`a-dire usuelle).
      Cette \'ecriture signifie que 
            \begin{align*}
            N & = a_{n} \cdot 10^{n} + a_{n-1} \cdot 10^{n-1} + \cdots 
                  + a_{2} \cdot 10^{2} + a_{1} \cdot 10^{1} + a_{0} \cdot 10^{0}            \\
              & = \sum_{i=0}^{n} a_{i} \cdot 10^{i}
            \end{align*}
      \begin{ex}
            Consid\'erons le nombre $N = 9\;235\;819$.
            On a alors 
            $N = 9\cdot 10^{6} + 2\cdot 10^{5} + 3\cdot 10^{4} + 5\cdot 10^{3} + 8\cdot 10^{2} + 1\cdot 10^{1} + 9\cdot 10^{0}$.
      \end{ex}
      
      Pour \'ecrire un nombre entier naturel $N$ en base $b$ (o\`u $b\geq 2$), on se sert de $b$ symboles.
      Si l'\'ecriture en base $b$ de $N$ est 
            $ a_n a_{n-1} \dots a_2 a_1 a_0 $
      alors 
            \begin{align*}
            N & = a_{n} \cdot b^{n} + a_{n-1} \cdot b^{n-1} + \cdots 
                  + a_{2} \cdot b^{2} + a_{1} \cdot b^{1} + a_{0} \cdot b^{0}            \\
              & = \sum_{i=0}^{n} a_{i} \cdot b^{i}                                       .
            \end{align*}
      Pour indiquer que l'\'ecriture est en base $b$, on note \defbox{$ \overline{a_n a_{n-1} \dots a_2 a_1 a_0}^{(b)}$}.
      \begin{conversionexample}
      \'Ecrivons les nombres 13 et 41 en base 2 (\'ecriture binaire).
\index{Numeration@Num\'eration!binaire}
\index{Numeration@Num\'eration!base 2}
            \par
            \TCDConversionPair{%
            \begin{tikzpicture}[x=1.5em,y=1.2em,>=stealth]
                  \draw (-1.35,-.37) node {$13$};
                  \draw [rounded corners = 4,color=loko2,->,line width=3,opacity=.3] 
                        (1.65,-3.25) -- (0.7,-4.45) -- (-1.3,-2.25);
                  \foreach \x in {0,1,2}
                        \draw (\x-.75,-\x-.1) -- ++(0,-.75) -- ++(.75,0) ++(-.35,.5) node{$2$};
                  \foreach \x/\y/\z in {0/6/12,1/3/6,2/\textcolor{loko1}{1}/2}
                        \draw (\x-.35,-\x-1.35) node{$\y$} (\x-1.35,-\x-1.32) node {$\z$};
                  \foreach \x/\y in {0/{\phantom{0}1},1/0,2/1}
                        \draw (\x-1.65,-\x-1.8) -- (\x-1,-\x-1.8) 
                              (\x-1.3,-\x-2.3) node{\textcolor{loko1}{$\y$}};
            \end{tikzpicture}
            }{Ainsi, $13 = \overline{1101}^{_{(2)}}$.}{%
            \begin{tikzpicture}[x=1.5em,y=1.2em,>=stealth]
                  \draw (-1.35,-.37) node {$41$};
                  \draw [rounded corners = 4,color=loko2,->,line width=3,opacity=.3] 
                        (3.65,-5.25) -- (2.7,-6.45) -- (-1.3,-2.25);
                  \foreach \x in {0,1,2,3,4}
                        \draw (\x-.75,-\x-.1) -- ++(0,-.75) -- ++(.75,0) ++(-.35,.5) node{$2$};
                  \foreach \x/\y/\z in {0/20/40,1/10/20,2/5/10,3/2/4,4/\textcolor{loko1}{1}/2}
                        \draw (\x-.35,-\x-1.35) node{$\y$} (\x-1.35,-\x-1.32) node {$\z$};
                  \foreach \x/\y in {0/{\phantom{0}1},1/0,2/0,3/1,4/0}
                        \draw (\x-1.65,-\x-1.8) -- (\x-1,-\x-1.8) 
                              (\x-1.3,-\x-2.3) node{\textcolor{loko1}{$\y$}};
            \end{tikzpicture}
            }{Ainsi, $41 = \overline{101001}^{_{(2)}}$.}
      \end{conversionexample}
      %
      \begin{conversionexample}
            \'Ecrivons les nombres 957 et 188899 en base 8 :
            \par
            \TCDConversionPair{%
            \begin{tikzpicture}[x=2em,y=1.1em,>=stealth]
                  \draw (-1.35,-.37) node {$957$};
                  \draw [rounded corners = 4,color=loko2,->,line width=3,opacity=.3] 
                        (1.65,-3.25) -- (0.7,-4.45) -- (-1.3,-2.25);
                  \foreach \x in {0,1,2}
                        \draw (\x-.75,-\x-.1) -- ++(0,-.75) -- ++(.75,0) ++(-.35,.5) node{$8$};
                  \foreach \x/\y/\z in {0/119/952,1/14/112,2/\textcolor{loko1}{1}/{\phantom{0}8}}
                        \draw (\x-.35,-\x-1.35) node{$\y$} (\x-1.35,-\x-1.32) node {$\z$};
                  \foreach \x/\y in {0/{\phantom{00}5},1/{\phantom{00}7},2/{\phantom{0}6}}
                        \draw (\x-1.65,-\x-1.8) -- (\x-1,-\x-1.8) 
                              (\x-1.35,-\x-2.3) node{\textcolor{loko1}{$\y$}};
            \end{tikzpicture}
            }{Ainsi, $957 = \overline{1675}^{_{(8)}}$.}{%
            \begin{tikzpicture}[x=3.5em,y=1.1em,>=stealth]
                  \draw (-1.35,-.37) node {$188899$};
                  \draw [rounded corners = 4,color=loko2,->,line width=3,opacity=.3] 
                        (3.65,-5.25) -- (2.7,-6.45) -- (-1.2,-2.25);
                  \foreach \x in {0,1,2,3,4}
                        \draw (\x-.75,-\x-.1) -- ++(0,-.75) -- ++(.75,0) ++(-.35,.5) node{$8$};
                  \foreach \x/\y/\z in {0/23612/188896,1/2951/23608,2/368/2944,3/46/368,4/\textcolor{loko1}{5}/40}
                        \draw (\x-.35,-\x-1.35) node{$\y$} (\x-1.35,-\x-1.32) node {\underline{$\z$}};
                  \foreach \x/\y in {0/{\phantom{0000}3},1/{\phantom{000}4},2/{\phantom{00}7},3/{\phantom{0}0},4/{6}}
                        \draw (\x-1.3,-\x-2.3) node{\textcolor{loko1}{$\y$}};
            \end{tikzpicture}
            }{Ainsi, $188899 = \overline{560743}^{_{(8)}}$.}
      \end{conversionexample}
      %
      \begin{conversionexample}
            \'Ecrivons les nombres 18 647 et 1 304 748 en base 12, pour cel\`a, 
            on aura besoin d'ajouter les symboles $A$ et $B$ pour les <<chiffres>> dix et onze respectivement.
            \par
            \TCDConversionPair{%
            \begin{tikzpicture}[x=3em,y=1.1em,>=stealth]
                  \draw (-1.35,-.37) node {$18 647$};
                  \draw [rounded corners = 4,color=loko2,->,line width=3,opacity=.3] 
                        (1.7,-3.25) -- (0.8,-4.45) -- (-1.2,-2.25);
                  \foreach \x in {0,1,2}
                        \draw (\x-.75,-\x-.1) -- ++(0,-.75) -- ++(.75,0) ++(-.35,.5) node{$12$};
                  \foreach \x/\y/\z in {0/1553/18636,1/129/1548,2/\textcolor{loko1}{10}/{120}}
                        \draw (\x-.35,-\x-1.35) node{$\y$} (\x-1.35,-\x-1.32) node {\underline{$\z$}};
                  \foreach \x/\y in {0/{\phantom{000}11},1/{\phantom{000}5},2/{\phantom{00}9}}
                        \draw (\x-1.35,-\x-2.3) node{\textcolor{loko1}{$\y$}};
            \end{tikzpicture}
            }{Ainsi, $18\,647 = \overline{A95B}^{_{(12)}}$.}{%
            \begin{tikzpicture}[x=4em,y=1.1em,>=stealth]
                  \draw (-1.35,-.37) node {$1 304 748$};
                  \draw [rounded corners = 4,color=loko2,->,line width=3,opacity=.3] 
                        (3.65,-5.25) -- (2.7,-6.45) -- (-1.2,-2.25);
                  \foreach \x in {0,1,2,3,4}
                        \draw (\x-.75,-\x-.1) -- ++(0,-.75) -- ++(.75,0) ++(-.35,.5) node{$12$};
                  \foreach \x/\y/\z in {0/108729/1304748,1/9060/108720,2/755/9060,3/62/744,4/\textcolor{loko1}{5}/60}
                        \draw (\x-.35,-\x-1.35) node{$\y$} (\x-1.35,-\x-1.32) node {\underline{$\z$}};
                  \foreach \x/\y in
                           {0/{\phantom{00000}0},1/{\phantom{0000}9},2/{\phantom{00}0},3/{\phantom{!}11},4/{\phantom{.}2}}
                        \draw (\x-1.3,-\x-2.3) node{\textcolor{loko1}{$\y$}};
            \end{tikzpicture}
            }{Ainsi, $1\,304\,748 = \overline{52B090}^{_{(12)}}$.}
      \end{conversionexample}
      %
      \begin{exe}\label{exe:arith-2}
            \'Ecrire les nombres suivants en base dix (num\'eration d\'ecimale) :
\index{Numeration@Num\'eration!decimale@d\'ecimale}
\index{Numeration@Num\'eration!base 10}
            $ \overline{101110}^{_{(2)}} $, $ \overline{756}^{_{(8)}} $, $ \overline{B00A16}^{_{(12)}} $.
      \end{exe}
      %
      \warning{%
            Il existe d'autres notations pour les nombres en base autres que dix.
            Le nombre $42$ peut par exemple \^etre not\'e
                  $ \overline{101010}^{_{(2)}} $,
                  $ \overline{101010}^{_{2}} $,
                  $ (\overline{101010})_{2} $,
                  $ {101010}_{2} $,
                  \ldots.
            On choisira, dans ce cours, d'utiliser toujours la notation $ \overline{101010}^{_{(2)}} $.
            \newline
            Pour les bases sup\'erieures \`a dix,
            on ajoute souvent des lettres de l'alphabet latin dans l'ordre.
            On peut utiliser des lettres minuscules aussi bien que capitales.
      }
\section{Congruence}
\index{Congruence}
      \begin{defi}
            Soient $(x,y) \in \mathbb{Z}^2$ et $n\in\mathbb{N}$.
            On dit que \emph{$x$ est congru \`a $y$ modulo $n$} si
            \emph{$n$ divise $x-y$}. 
            On note \defbox{$x \equiv y \quad \mathrm{mod}\ n$} ou \defbox{$x \equiv y \ [n]$}.
            %On note $\congru{x}{y}{n}$.
                  \begin{IEEEeqnarray*}{rCl}
                        \congru{x}{y}{n} & \iff & n | x - y                                     \\
                                         & \iff & x - y = k \cdot n \quad (k\in\mathbb{Z})      \\
                                         & \iff & x = y + k \cdot n \quad (k\in\mathbb{Z})      
                  \end{IEEEeqnarray*}
            \smallnota{La congruence modulo $0$ est l'\'egalit\'e.}
      \end{defi}
      %
      \begin{thm}
            Soit $n$ un entier naturel.
            \begin{enumerate}
                  \item La congruence modulo $n$ est une relation d'\'equivalence, 
                        c'est-\`a-dire que, pour tous entiers relatifs $x$, $y$ et $z$, on a :
\index{Congruence!relation d'equivalence@relation d'\'equivalence}
                        \begin{description}
                              \item[Reflexivit\'e :] 
                                    $\congru{x}{x}{n}$,
                              \item[Sym\'etrie :] 
                                    $\congru{x}{y}{n} \iff \congru{y}{x}{n}$,
                              \item[Transitivit\'e :] 
                                    $\left( \congru{x}{y}{n} \text{ et } \congru{y}{z}{n} \right) \Rightarrow \congru{x}{z}{n}$.
                        \end{description}
                  \item La congruence modulo $n$ est compatible avec l'addition et la multiplication, 
                        c'est-\`a-dire que, pour tous entiers relatifs $x_1$, $x_2$, $y_1$ et $y_2$, on a :
\index{Congruence!compatibilite avec l'addition@compatibilit\'e avec l'addition}
\index{Congruence!compatibilite avec la multiplication@compatibilit\'e avec la multiplication}
                        \begin{displaymath}
                              \left.
                                    \begin{array}{l}
                                          \congru{x_1}{y_1}{n} \\
                                          \congru{x_2}{y_2}{n}
                                    \end{array}
                              \right\rbrace
                              \Rightarrow
                              \left\lbrace
                                    \begin{array}{l}
                                          \congru{x_1 + x_2}{y_1 + y_2}{n}             \\
                                          \congru{x_1 \cdot x_2}{y_1 \cdot y_2}{n}     \\
                                    \end{array}
                              \right.
                        \end{displaymath}
                        et par cons\'equent, pour tout $k\in\mathbb{N}^*$, on a :
                              $\congru{x}{y}{n} \Rightarrow \congru{x^k}{y^k}{n}$.
            \end{enumerate}
      \end{thm}
      %
      \begin{prop}
            \begin{enumerate}
                  \item Si $d|n$ et $\congru{x}{y}{n}$, alors $\congru{x}{y}{d}$.
                  \item Soit $n\in\mathbb{N}^*$.
                        Pour tout $x\in\mathbb{Z}$, il existe un unique $r\in\mathbb{N}$ tel que
                        $0\leq r < n$ et $\congru{x}{r}{n}$.
            \end{enumerate}
      \end{prop}
      %
      \begin{ex}
            On veut d\'eterminer le reste de la division euclidienne de $7^{100}$ par $5$.
            
            On a $7^{100} = (7^2)^{50} = 49^{50}$.
            Comme $49 = 50 - 1 = 5 \cdot 10 - 1$,
            alors  $\congru{49}{-1}{5}$ ; c'est-\`a-dire $\congru{7^2}{-1}{5}$.
            Il vient que $\congru{(7^2)^{50}}{(-1)^{50}}{5}$ et donc $\congru{7^{100}}{1}{5}$.
            Ainsi, le reste de la division euclidienne de $7^{100}$ par $5$ est $1$.
      \end{ex}
      %
      \begin{exe}\label{exe:arith-3}
            \begin{enumerate}
                  \item D\'eterminer le reste de la division euclidienne de $5^{2018}$ par $13$.
                  \item D\'eterminer le reste de la division euclidienne de $2^{2019}$ par $7$.
                  \item D\'eterminer le reste de la division euclidienne de $5^{4037}$ par $7$.
            \end{enumerate}
      \end{exe}
      %
      \begin{ex}
\index{Critere de divisibilite par 9@Crit\`ere de divisibilit\'e par $9$}
            On veut d\'eterminer un crit\`ere de divisibilit\'e par $9$.
            
            Soit $N$ un nombre entier naturel et soit $a_n a_{n-1} \ldots a_2 a_1 a_0$ son \'ecriture en base 10.
            On a alors $N = a_n \cdot 10^n + a_{n-1} \cdot 10^{n-1} + \cdots + a_1 \cdot 10 + a_0$.
            
            Montrons, d'abord, que pour tout $p\in\mathbb{N}$, on a $\congru{10^p}{1}{9}$ :
                  \begin{align*}
                        10 & = 9 \cdot 1 + 1     \\
                        \congru{10 & }{1}{9}     \\
                        \congru{10^p & }{1^p}{9} \\
                        \congru{10^p & }{1}{9}   \;.
                  \end{align*}
                  
            Ainsi, pour tout $p\in\mathbb{N}$, on a :
                  \begin{align*}
                        \congru{10^p & }{1}{9}                   \\
                        \congru{a_p \cdot 10^p & }{a_p}{9}       \;.
                  \end{align*}
            En prenant la somme pour $p$ allant de $0$ \`a $n$, on a :
                  \begin{align*}
                        \congru{a_n10^n+a_{n-1}10^{n-1}+\cdots+a_110+a_0}{a_n+a_{n-1}+\cdots+a_1+a_0}{9}      \;.
                  \end{align*}
            Ainsi, $\congru{N}{\sum_{p=0}^n a_p}{9}$.
            Comme $9|N \iff \congru{N}{0}{9}$, alors $9|N \iff \congru{\sum_{p=0}^n a_p}{0}{9}$.
            Ainsi, $N$ est un multiple de $9$ si et seulement si la somme de ces chiffres est un multiple de $9$.
      \end{ex}
      %
      \begin{exe}\label{exe:arith-4}
            Soit $N$ un nombre entier naturel et soit $a_n a_{n-1} \ldots a_2 a_1 a_0$ son \'ecriture en base 10.
            \begin{enumerate}
                  \item Montrer que $N$ est divisible par $2$ si et seulement si $\congru{a_0}{0}{2}$.
                  \item Montrer que $N$ est divisible par $5$ si et seulement si $\congru{a_0}{0}{5}$.
                  \item Montrer que $3|N \iff \congru{\sum_{p=0}^n a_p}{0}{3}$.
            \end{enumerate}
      \end{exe}
      %
      \section{Anneau $\mathbb{Z}/n\mathbb{Z}$}
\index{Anneau}
\index{Z/nZ@$\mathbb{Z}/n\mathbb{Z}$}
      Soient $n\in\mathbb{N}$, $n \geq 2$, et $r\in\mathbb{Z}$.
      On d\'efinit la classe d'\'equivalence de $r$ comme \'etant 
      l'ensemble des $x\in\mathbb{Z}$ tel que $\congru{x}{r}{n}$, 
      on la note par \defbox{$\bar{r}$} ou \defbox{$\dot{r}$}.
      Ainsi,
      \defbox{$\bar{r} = \left\lbrace x\in\mathbb{Z} \mid \congru{x}{r}{n} \right\rbrace$}.
\index{Classe d'equivalence@Classe d'\'equivalence}
      %
      Bien qu'une classe d'\'equivalence peut avoir plusieurs repr\'esentants, 
      on pr\'ef\`ere prendre $ r \in \llbracket 0 ; n-1 \rrbracket$.
      %
      On note par $\mathbb{Z}/n\mathbb{Z}$ l'ensemble 
      $\left\lbrace \bar{0},\ \bar{1},\ \ldots ,\ \bar{n-1} \right\rbrace$.
      %
      On d\'efinit l'addition et la multiplication dans $\mathbb{Z}/n\mathbb{Z}$ par :       
      \begin{align*}
            \bar{a} + \bar{b} & = \bar{a+b}       & \bar{a} \cdot \bar{b} = \bar{a\cdot b} \;.
      \end{align*}
      Ces op\'erations ne d\'ependent pas des repr\'esentants choisis (compatibilit\'e de la congruence avec l'addition et la multiplication).
      Muni de ces deux op\'erations, $\mathbb{Z}/n\mathbb{Z}$ a une structure d'anneau,
\index{Anneau}
      c'est-\`a-dire que pour tout $\bar{a}$, $\bar{b}$ et $\bar{c}$ de $\mathbb{Z}/n\mathbb{Z}$, on a :
            \begin{itemize}
                  \item $\bar{a} + \bar{b}$ et $\bar{a} \cdot \bar{b}$ sont bien d\'efinie dans $\mathbb{Z}/n\mathbb{Z}$,
                  \item $(\bar{a} + \bar{b}) + \bar{c} = \bar{a} + (\bar{b} + \bar{c})$,
                  \item $\bar{a} + \bar{b} = \bar{b} + \bar{a}$,
                  \item $\bar{a} + \bar{0} = \bar{a}$,
                  \item il existe un \'el\'ement, not\'e $-\bar{a}$ tel que $\bar{a} + (-\bar{a}) = \bar{0}$,
                  \item $(\bar{a} \cdot \bar{b}) \cdot \bar{c} = \bar{a} \cdot (\bar{b} \cdot \bar{c})$,
                  \item $\bar{a} \cdot (\bar{b} + \bar{c}) = \bar{a} \cdot \bar{b} + \bar{a} \cdot \bar{c}$,
                  \item $(\bar{b} + \bar{c}) \cdot \bar{a} = \bar{b} \cdot \bar{a} + \bar{c} \cdot \bar{a}$.
            \end{itemize}
      %
      \begin{ex}
            $\mathbb{Z}/5\mathbb{Z} = \left\lbrace \bar{0},\ \bar{1},\ \bar{2},\ \bar{3},\ \bar{4} \right\rbrace$.
            Dans ce cas, 
            \begin{align*}
                  \bar{0} & = \left\lbrace \ldots ,\ -15,\ -10,\ -5,\ 0,\ 5,\ 10,\ 15,\ \ldots \right\rbrace ,\\
                  \bar{1} & = \left\lbrace \ldots ,\ -14,\ -9,\ -4,\ 1,\ 6,\ 11,\ 16,\ \ldots \right\rbrace  ,\\
                  \bar{2} & = \left\lbrace \ldots ,\ -13,\ -8,\ -3,\ 2,\ 7,\ 12,\ 17,\ \ldots \right\rbrace  ,\\
                  \bar{3} & = \left\lbrace \ldots ,\ -12,\ -7,\ -2,\ 3,\ 8,\ 13,\ 18,\ \ldots \right\rbrace  ,\\
                  \bar{4} & = \left\lbrace \ldots ,\ -11,\ -6,\ -1,\ 4,\ 9,\ 14,\ 19,\ \ldots \right\rbrace  .
            \end{align*}
            On a les tables d'additions et de multiplications suivantes :
                  \[
                  \begin{array}{|c|c|c|c|c|c|}
                  \hline 
                  + & \bar{0} & \bar{1} & \bar{2} & \bar{3} & \bar{4} \\ 
                  \hline 
                  \bar{0} & \bar{0} & \bar{1} & \bar{2} & \bar{3} & \bar{4} \\ 
                  \hline 
                  \bar{1} & \bar{1} & \bar{2} & \bar{3} & \bar{4} & \bar{0} \\ 
                  \hline 
                  \bar{2} & \bar{2} & \bar{3} & \bar{4} & \bar{0} & \bar{1} \\ 
                  \hline 
                  \bar{3} & \bar{3} & \bar{4} & \bar{0} & \bar{1} & \bar{2} \\ 
                  \hline 
                  \bar{4} & \bar{4} & \bar{0} & \bar{1} & \bar{2} & \bar{3} \\ 
                  \hline 
                  \end{array}
                  \qquad
                  \begin{array}{|c|c|c|c|c|c|}
                  \hline 
                  \times & \bar{0} & \bar{1} & \bar{2} & \bar{3} & \bar{4} \\ 
                  \hline 
                  \bar{0} & \bar{0} & \bar{0} & \bar{0} & \bar{0} & \bar{0} \\ 
                  \hline 
                  \bar{1} & \bar{0} & \bar{1} & \bar{2} & \bar{3} & \bar{4} \\ 
                  \hline 
                  \bar{2} & \bar{0} & \bar{2} & \bar{4} & \bar{1} & \bar{3} \\ 
                  \hline 
                  \bar{3} & \bar{0} & \bar{3} & \bar{1} & \bar{4} & \bar{2} \\ 
                  \hline 
                  \bar{4} & \bar{0} & \bar{4} & \bar{3} & \bar{2} & \bar{1} \\ 
                  \hline 
                  \end{array}
                  .
                  \] 
      \end{ex}
      %
      \smallnota{Au lieu de $\mathbb{Z}/n\mathbb{Z}$, on utilise parfois $\mathbb{Z}_n$.}
      %
      \begin{ex}                  
            R\'esolvons dans $\mathbb{Z}/5\mathbb{Z}$ l'\'equation $x+\bar{3} = \bar{2}$.
            Pour cel\`a, additionnons $\bar{2}$ dans chaque membre de l'\'egalit\'e, on a alors 
                  $x + \bar{3} + \bar{2} = \bar{2} + \bar{2} \Rightarrow x = \bar{4}$.
            On aurait aussi pu utiliser la m\'ethode usuelle et trouver $ x = \bar{-1} = \bar{4}$.
            
            R\'esolvons dans $\mathbb{Z}/5\mathbb{Z}$ l'\'equation $\bar{3}\, x+\bar{3} = \bar{2}$.
            On sait d\'ej\`a (en utilisant ce qui pr\'ec\`ede) que $\bar{3} x = \bar{4}$.
            En multipliant l'\'egalit\'e par $\bar{2}$, on obtient 
                  $\bar{2} \cdot \bar{3} \cdot x = \bar{2} \cdot \bar{4}$
                  et donc $x = \bar{3}$.
            
            R\'esolvons dans $\mathbb{Z}/5\mathbb{Z} \times \mathbb{Z}/5\mathbb{Z}$ le syst\`eme d'\'equations :
            \[
            \left\lbrace
                  \begin{array}{l}
                        \bar{3}\, x + \bar{2}\, y = \bar{1} \\
                        \bar{2}\, x + \bar{4}\, y = \bar{0} \;.
                  \end{array}
            \right.
            \]
            On peut additionner les deux \'equations pour avoir 
                $ ( \bar{3} + \bar{2} ) \, x + ( \bar{2} + \bar{4} ) \, y = \bar{1} + \bar{0} $.
                Ainsi, $y = \bar{1}$.
            Pour trouver la valeur de $x$, on peut commencer par multiplier la premi\`ere \'equation par $\bar{3}$,
            \[
            \left\lbrace
                  \begin{array}{l}
                        \bar{4}\, x + \bar{1}\, y = \bar{3} \\
                        \bar{2}\, x + \bar{4}\, y = \bar{0} \\
                  \end{array}
            \right.
            \]
            et ensuite additionner les deux \'equations pour avoir 
                  $ x = \bar{3}$.            
      \end{ex}
      %
      \begin{exe}\label{exe:arith-5}
            \begin{enumerate}
            \item R\'esoudre dans $\mathbb{Z}/7\mathbb{Z}$ les \'equations suivantes : 
                  \begin{enumerate}
                        \item $\bar{3}\, x + \bar{4} = \bar{5}$,
                        \item $\bar{2}\, x + \bar{6} = \bar{3}$,
                        \item $\bar{6}\, x + \bar{5} = \bar{1}$.
                  \end{enumerate}
            \item R\'esoudre dans
                  $\mathbb{Z}/7\mathbb{Z}\times \mathbb{Z}/7\mathbb{Z}$
                  les syst\`emes d'\'equations suivants :
                  \begin{enumerate}
                        \item
                              $\left\lbrace
                                    \begin{array}{l}
                                          \bar{2}x + \bar{5}y = \bar{3} \\
                                          \bar{3}x + \bar{2}y = \bar{1} \,;
                                    \end{array}
                               \right.$
                        \item
                              $\left\lbrace
                                    \begin{array}{l}
                                                  x + \bar{2}y = \bar{6} \\
                                          \bar{2}x + \bar{3}y = \bar{4} \,;
                                    \end{array}
                               \right.$
                        \item
                              $\left\lbrace
                                    \begin{array}{l}
                                          \bar{5}x +         y = \bar{3} \\
                                          \bar{3}x + \bar{2}y = \bar{0} \,;
                                    \end{array}
                               \right.$
                        \item
                              $\left\lbrace
                                    \begin{array}{l}
                                          \bar{2}x + \bar{3}y = \bar{1} \\
                                          \bar{3}x +         y = \bar{5} \,.
                                    \end{array}
                               \right.$
                  \end{enumerate}
            \end{enumerate}
      \end{exe}
      %
\section{PPCM et PGCD}
      \subsection{PPCM}
      \begin{defi}
\index{PPCM}
      Soient $a$ et $b$ deux entiers naturels non nuls.
      On appelle \textbf{P}lus \textbf{P}etit \textbf{C}ommun \textbf{M}ultiple 
      --- en abr\'eg\'e \emph{PPCM} --- de $a$ et $b$ le nombre 
            \defbox{$PPCM(a;b) := \min \left( a\mathbb{N}^* \cap b\mathbb{N}^* \right)$}
            o\`u $x\mathbb{N}^* := \left\lbrace x \cdot n \mid n \in \mathbb{N}^* \right\rbrace$.
      
      \smallnota{%
                  Le PPCM de $a$ et $b$ est aussi parfois not\'e $a \vee b$ ou $[a;b]$.
                 }
                 
      Si $a$ et $b$ sont deux entiers relatifs, 
      on peut d\'efinir le PPCM de $a$ et $b$ comme \'etant le PPCM de $|a|$ et $|b|$.
      Si $a$ et $b$ sont deux entiers relatifs non-nuls, on a aussi :
            \defbox{$PPCM(a;b) := \min \left( \left\lbrace m\in\mathbb{N}^* \mid a|m \text{ et } b|m \right\rbrace \right)$.}
      \end{defi}
      %
      \begin{prop}
            Soient $a$ et $b$ deux entiers naturels non-nuls.
            \begin{itemize}
                  \item Si $a|b$ alors $PPCM(a;b) = b$,
                  \item si $k\in\mathbb{N}^*$ alors $PPCM(ka;kb) = k \cdot PPCM(a;b)$,
                  \item si $c\in\mathbb{N}^*$, $PPCM(a;b;c) = PPCM(PPCM(a;b);c) = PPCM(a;PPCM(b;c))$
                        (on peut \'etendre \`a un nombre arbitraire d'\'el\'ements).
            \end{itemize}
      \end{prop}
      %
      \subsection*{Recherche du PPCM}
            Pour trouver le PPCM de deux --- ou de plusieurs --- nombres, 
            on peut d\'ecomposer chacun des nombres en produits de facteurs premiers.
\index{Decomposition en produits de facteurs premiers@D\'ecomposition en produits de facteurs premiers}
            Il suffit ensuite de prendre, pour chaque facteur premier pr\'esent, la plus grande puissance.
            
            \begin{ex}
                  Cherchons $PPCM(28;24;36)$.
                  
                  Comme 
                  $
                  \left\lbrace
                  \begin{array}{l}
                        28 = 2^2 \cdot 7 \\
                        24 = 2^3 \cdot 3 \\
                        36 = 2^2 \cdot 3^2 \\
                  \end{array}
                  \right.
                  $
                  le PPCM de ces trois nombres est $2^3 \cdot 3^2 \cdot 7 = 504$.
            \end{ex}

      \subsection{PGCD}
      \begin{defi}
\index{PGCD}
      Soient $a$ et $b$ deux entiers naturels non tous les deux nuls.
      On appelle \textbf{P}lus \textbf{G}rand \textbf{C}ommun \textbf{D}iviseur
      --- en abr\'eg\'e \emph{PGCD} --- de $a$ et $b$ le nombre 
            \defbox{$\pgcd{a}{b} := \max \left( D_a \cap D_b \right)$}.
      
      \nota{%
                  Le PGCD de $a$ et $b$ est aussi parfois not\'e $\pgcd{a}{b}$, $a \wedge b$ ou $(a;b)$. \newline
                  On trouve parfois l'acronyme PGDC 
                  (\textbf{P}lus \textbf{G}rand \textbf{D}iviseur \textbf{C}ommun)
\index{PGDC}
                  mais la version la plus r\'epandue est PGCD.
                 }
      
      Deux nombres $a$ et $b$ sont dits \emph{premiers entre eux} si $\pgcd{a}{b} = 1$.
\index{Premiers entre eux}
      \end{defi}
      \begin{rmk}
            Tout entier divise $0$, donc $D_0 \cap D_0$ n'a pas de plus grand \'el\'ement ;
            on pose, par convention, $\pgcd{0}{0}=0$.
            Pour $a \neq 0$, on a $\pgcd{a}{0} = a$.
            Pour des entiers relatifs, on pose $\pgcd{a}{b} = \pgcd{|a|}{|b|}$.
      \end{rmk}
      %
      \begin{prop}
            Soient $a$ et $b$ deux entiers naturels non-nuls.
            \begin{itemize}
                  \item Si $a|b$ alors $\pgcd{a}{b} = a$,
                  \item si $k\in\mathbb{N}^*$ alors $\pgcd{ka}{kb} = k \cdot \pgcd{a}{b}$,
                  \item si $c\in\mathbb{N}^*$, $\mathrm{PGCD}(a,b,c) = \pgcd{\pgcd{a}{b}}{c} = \pgcd{a}{\pgcd{b}{c}}$
                        (on peut \'etendre \`a un nombre arbitraire d'\'el\'ements).
            \end{itemize}
      \end{prop}
      %
      \begin{thm}
            Soient $a$ et $b$ deux nombres entiers naturels non-nuls.
            Soit $r$ le reste de la division de $b$ par $a$.
            %
            Alors $\pgcd{a}{b}$ divise $r$.
            
            \smallnota{En fait, on a un r\'esultat plus fort : \thmbox{$\pgcd{a}{b} = \pgcd{a}{r}$}
                       (car $r = b - qa$ et $b = qa + r$ : les diviseurs communs de $a$ et $b$ sont exactement ceux de $a$ et $r$).}
      \end{thm}
      %
      \begin{thm}
            Soient $a$ et $b$ deux entiers naturels non-nuls.
            On a
            \[
                  \thmbox{PPCM(a;b)\cdot\pgcd{a}{b}=a\cdot b}.
            \]
      \end{thm}
      %
      \subsection*{Recherche du PGCD}
            Pour trouver le PGCD de deux --- ou de plusieurs --- nombres, 
            on peut d\'ecomposer chacun des nombres en produits de facteurs premiers.
\index{Decomposition en produits de facteurs premiers@D\'ecomposition en produits de facteurs premiers}
            Il suffit ensuite de prendre, pour chaque facteur premier pr\'esent, la plus petite puissance
            (\'eventuellement $0$).
            
            \begin{ex}
                  Cherchons $\pgcd{36}{56}$.
                  
                  Comme 
                  $
                  \left\lbrace
                  \begin{array}{l}
                        36 = 2^2 \cdot 3^2 \\
                        56 = 2^3 \cdot 7 \\
                  \end{array}
                  \right.
                  $
                  le PGCD de ces deux nombres est $2^2 = 4$.
            \end{ex}
      \subsection*{Algorithme d'Euclide}
\index{Algorithme d'Euclide}
\index{Euclide!Algorithme de}
            Il existe un algorithme --- celui d'Euclide --- permettant de trouver le PGCD de deux nombres entiers naturels.
            
            Soient $a$ et $b$ deux entiers naturels non-nuls tels que $a>b$.
            On pose $a_0=a$, $a_1=b$ et, pour tout entier naturel $n$,
            \[
              \begin{aligned}
                a_n &= q_n a_{n+1}+r_n,
                & 0 &\leq r_n<a_{n+1}, \\
                a_{n+2} &= r_n.
              \end{aligned}
            \]
            Les nombres $q_n$ et $r_n$ sont respectivement le quotient et le
            reste de la division de $a_n$ par $a_{n+1}$.
                  
            On peut dresser le tableau suivant
            \[
            \begin{array}{|c|c|c|c|c|}
            \hline 
              & q_0 & q_1 & q_2 & q_3 \\ 
            \hline 
            a_0 = a & a_1 = b & a_2 = r_0 & a_3 = r_1  & \cdots \\ 
            \hline 
            r_0 = r & r_1 = a_1 - a_2 \cdot q_1 & r_2 = a_2 - a_3 \cdot q_2 & \cdots & \cdots \\ 
            \hline 
            \end{array} 
            \]
            et le dernier reste (c'est dans la derni\`ere ligne) non-nul est le PGCD de $a$ et $b$.
            
            \begin{ex}
            \label{ex:Algorithme d'Euclide: 36,200}
                  Cherchons le PGCD de $36$ et $200$.
                  \[
                  \begin{array}{|c|c|c|c|c|}
                  \hline 
                           & q_0 = 5  & q_1 = 1   & q_2 = 1 & q_3 = 4 \\ 
                  \hline 
                  200      &      36  &      20   &      16 &       4 \\ 
                  \hline 
                  r_0 = 20 & r_1 = 16 & r_2 = 4   &       0 &         \\ 
                  \hline 
                  \end{array} 
                  \]
                  Le dernier reste non-nul est $4$, donc $\pgcd{36}{200} = 4$.
            \end{ex}
            %
            \begin{exe}\label{exe:arith-6}
                  Calculer :
                  \begin{enumerate}[label=\alph*)]
                        \item $\pgcd{51\,959}{47\,897}$ ;
                        \item $\pgcd{23\,213}{24\,881}$ ;
                        \item $\pgcd{47\,027}{62\,651}$.
                  \end{enumerate}
            \end{exe}
            %
            \begin{thm}[Identit\'e de B\'ezout]
\index{Bezout@B\'ezout!Identite de@Identit\'e de}
%\index{Bezout@B\'ezout!Theoreme de@Th\'eor\`eme de}
%\index{Bezout@B\'ezout!Egalite de@\'Egalit\'e de}
                  Soient $a$ et $b$ deux entiers relatifs non-nuls et notons par $d$ le PGCD de $a$ et $b$.
                  On a alors :
                  \begin{itemize}
                        \item Il existe deux entiers $u$ et $v$ v\'erifiant : $a u + b v = d$,
                        \item pour tous entiers $u$ et $v$, le nombre $a u + b v$ est un multiple de $d$.
                  \end{itemize}
                  Les nombres $u$ et $v$ sont appel\'es \emph{coefficients de B\'ezout},
                  ils ne sont pas uniques.
\index{Bezout@B\'ezout!coefficients de}
\index{coefficients de B\'ezout@coefficients de Bezout}
            \end{thm}
            %
%            \begin{thm}[B\'ezout]
%                  Soient $a$ et $b$ deux nombre entiers relatifs non-nuls.
%                  Alors il existe $(u,v) \in \mathbb{Z} \times \mathbb{Z}$ tels que 
%                  \thmbox{$au+bv=\pgcd{a}{b}$}.
%            \end{thm}
            %
            \begin{ex}
                  On sait que $\pgcd{200}{36} = 4$.
                  Cherchons deux entiers relatifs $u$ et $v$ tels que $200 u + 36 v = 4$.
                  
                  \begin{displaymath}
                        \begin{array}{|r@{\ =\ }l|r@{\ =\ }l|r@{\ =\ }l|}
                              \hline
                              \multicolumn{2}{|c|}{\text{Division}} &
	                                     \multicolumn{2}{|c|}{\text{\'Ecriture num\'erique} } & 
	                                           \multicolumn{2}{|c|}{\text{\'Ecriture litt\'erale} } \\
                              \hline
                              200 & 5 \cdot 36 + 20 & 20 & 200 - 5 \cdot 36       & r_0  & a   - 5 \cdot b    \\
                              \hline
                              36  & 1 \cdot 20 + 16 & 16 & 36  - 1 \cdot 20       & r_1  & b   - 1 \cdot r_0  \\
                              \hline
                              20  & 1 \cdot 16 + 4  & 4  & 20  - 1 \cdot 16       & r_2  & r_0 - 1 \cdot r_1  \\
                              \hline
                              16  & 4 \cdot 4  + 0  & \multicolumn{2}{|c|}{}      & \multicolumn{2}{|c|}{}  \\
                              \hline
                        \end{array}
                  \end{displaymath}
                  Ainsi, on a 
                        \begin{align*}
                              r_2 & = (a-5b) - (b-r_0)       \\
                                  & = (a-5b) - (b-(a-5b))    \\
                                  & = a-5b-b+a-5b            \\
                                  & = 2a - 11 b              \;.
                        \end{align*}
                 D'o\`u $4 = 2 \cdot 200 - 11 \cdot 36$.
            \end{ex}
            %
            \begin{exe}\label{exe:arith-7}
                  Pour chacun des couples $(a,b)$ suivants trouver $\pgcd{a}{b}$ :
                  %
                  \begin{enumerate}
                        \item $(a;b) = (615\,627 ; 36\,183)$, \label{item:1st gcd}
                        \item $(a;b) = (210\,573 ; 9\,099)$, \label{item:2nd gcd}
                        \item $(a;b) = (12\,325 ; 4\,207)$,
                        \item $(a;b) = (13\,567 ; 23\,369)$.
                  \end{enumerate}
                  %
                  Pour \ref{item:1st gcd} et \ref{item:2nd gcd}, trouver deux entiers $u$ et $v$ tels que
                  $au + bv = \pgcd{a}{b}$.
            \end{exe}
%
		\subsection*{Algorithme d'Euclide \'etendu}
\index{Algorithme d'Euclide!etendu@\'etendu}
\index{Euclide!Algorithme etendue d'@Algorithme \'etendu d'}
		Il existe une version ``\emph{\'etendue}'' de l'algorithme d'Euclide.
		Elle permet, en m\^eme temps que de conna\^itre le PGCD de deux nombres $a$ et $b$, 
		de trouver les coefficients de B\'ezout $u$ et $v$.
		
		Il est identique au pr\'ec\'edent, il faut juste ajouter deux lignes.
		\[
			\begin{array}{|c|c|c|c|c|}
				\hline 
					& q_0 & q_1 & q_2 & \cdots \\ 
				\hline 
					a_0 = a & a_1 = b & a_2 = r_0 & a_3 = r_1  & \cdots \\ 
				\hline 
					r_0 = r & r_1 = a_1 - a_2 \cdot q_1 & r_2 = a_2 - a_3 \cdot q_2 & r_3 = a_3 - a_4 \cdot q_3 & \cdots \\ 
				\hline
					\boldsymbol{u_0 = 1} & \boldsymbol{u_1 = 0} & u_2 = u_0 - (q_0 \cdot u_1) & u_3 = u_1 - (q_1 \cdot u_2) & \cdots\\
				\hline
					\boldsymbol{v_0 = 0} & \boldsymbol{v_1 = 1} & v_2 = v_0 - (q_0 \cdot v_1) & v_3 = v_1 - (q_1 \cdot v_2) & \cdots\\
				\hline
			\end{array} 
		\]
		Les derniers \'el\'ements des deux derni\`eres lignes sont les coefficients $u$ et $v$ qu'on cherche.
		
		En fait, \`a chaque \'etape, on a
		$
			r_{n-2} = a_{n} = u_{n} \; a + v_{n} \; b
		$.
		
		\begin{ex}
		Revenons \`a l'exemple pr\'ec\'edent.%~\ref{ex:Algorithme d'Euclide: 36,200}.
		On sait que $\pgcd{200}{36} = 4$.
		On veut utiliser l'algorithme d'Euclide pour trouver les nombres $u$ et $v$ tels que 
		$200 \, u + 36 \, v = 4$.
		
			\begin{tabular}{|r|r|r|r|r|}			\hline
					& $5$	& $1$	& $1$	& $4$		\\ \hline
			$200$	& $36$	& $20$	& $16$	& $4$		\\ \hline
			$20$	& $16$	& $4$	& $0$	&			\\ \hline
			$1$		& $0$	& $1$	& $-1$	& $2$		\\ \hline
			$0$		& $1$	& $-5$	& $6$	& $-11$	\\ \hline
			\end{tabular}
		
		Ainsi, on retrouve bien $200 \cdot 2 - 36 \cdot 11 = 4$.
		\end{ex}
            %
            \begin{thm}[Lemme de Gauss]
                  Soient $a$, $b$ et $c$ trois entiers relatifs tels que $\pgcd{a}{c} = 1$.
                  Alors, \thmbox{$a|bc \Rightarrow a | b$}.
            \end{thm}
            %
            \begin{ex}
                  On se propose de r\'esoudre dans $\mathbb{Z}\times\mathbb{Z}$ l'\'equation : 
                  $35 x + 38 y = 1$.
                  
                  Commen\c{c}ons par dresser l'algorithme d'Euclide pour trouver $\pgcd{38}{35}$ : \ 
                  \begin{tabular}{|l|l|l|l|l|} 
                  \hline
                             & $1$  & $11$ & $1$ & $2$ \\ 
                  \hline
                        $38$ & $35$ & $3$  & $2$ & $1$ \\ 
                  \hline
                        $3$  & $2$  & $1$  & $0$ &     \\
                  \hline
                  \end{tabular} .
                  
                  On peut alors chercher deux entiers $u$ et $v$ tels que $38 u + 35 v = 1$.
                                    \begin{displaymath}
                        \begin{array}{|r@{\ =\ }l|r@{\ =\ }l|r@{\ =\ }l|}
                              \hline
                              \multicolumn{2}{|c|}{\text{Division}} &
                                     \multicolumn{2}{|c|}{\text{Ecriture num\'erique} } & 
                                           \multicolumn{2}{|c|}{\text{\'Ecriture litt\'erale} } \\
                              \hline
                              38  & 1 \cdot 35 + 3 & 3 & 38 - 1 \cdot 35       & r_0  & a   - 1 \cdot b    \\
                              \hline
                              35  & 11 \cdot 3 + 2 & 2 & 35  - 11 \cdot 3      & r_1  & b   - 11 \cdot r_0 \\
                              \hline
                              3   & 1 \cdot 2 + 1  & 1  & 3  - 1 \cdot 2       & r_2  & r_0 - 1 \cdot r_1  \\
                              \hline
                              2  & 2 \cdot 1  + 0  & \multicolumn{2}{|c|}{}      & \multicolumn{2}{|c|}{}  \\
                              \hline
                        \end{array}
                  \end{displaymath}
                  Ainsi, on a : 
                        \begin{align*}
                              r_2 & = (a - b ) - (b - 11 \cdot r_0) \\
                                  & = a - b - b + 11 (a - b)        \\
                                  & = a - 2 b + 11 a - 11 b         \\
                                  & = 12 a - 13 b                   \;.
                        \end{align*}
                  
                  On aurait aussi pu utiliser l'algorithme d'Euclide \'etendu :
                  \begin{center}
							\begin{tabular}{|r|r|r|r|r|} \hline
							  & $1$ & $11$ & $1$ & $2$ \\ \hline
							$38$ & $35$ & $3$ & $2$ & $1$ \\ \hline
							$3$ & $2$ & $1$ & $0$ &   \\ \hline
							$1$ & $0$ & $1$ & $-11$ & $12$ \\ \hline
							$0$ & $1$ & $-1$ & $12$ & $-13$ \\ \hline
							\end{tabular}
                  \end{center}
                  
                  On a donc une solution particuli\`ere : $12 \cdot 38 - 13 \cdot 35 = 1$.
                  Si $u$ et $v$ sont deux solutions, on a : 
                        \begin{align*}
                        &
                              \left\lbrace
                                    \begin{array}{l}
                                          38 \cdot  u + 35 \cdot v  = 1 \\
                                          38 \cdot 12 - 35 \cdot 13 = 1 \\
                                    \end{array}
                              \right.
                        \Rightarrow
                        &     
                                    38 (u - 12) + 35 (v + 13) = 0      \;.
                        \end{align*}
                        Ainsi,
                              \begin{equation}
                                    35 \, (v + 13) = 38 (12 - u) \label{eq:Bezout ex 1}
                              \end{equation}
                        et donc $38 | 35 \, (v + 13)$.
                        Mais comme $\pgcd{38}{35} = 1$, d'apr\`es le Lemme~de~Gauss, $38 | (v + 13)$.
                        D'o\`u, il existe un entier $k$ tel que $v + 13 = 38 \, k$, ou encore $v = 38\,k - 13$.
                        En portant cette valeur dans l'\'equation~\eqref{eq:Bezout ex 1}, on a 
                              $38 \, (12 - u) = 35 \cdot 38 \cdot k$ et donc 
                              $u = - 35 \, k + 12$.
                        R\'eciproquement, pour tout $k\in\mathbb{Z}$, le couple $(u ; v) = (12 - 35\,k ; 38\,k - 13)$ v\'erifie
                        $38\,(12-35\,k) + 35\,(38\,k-13) = 456 - 455 = 1$.
                        Comme $x = v$ et $y = u$ (l'\'equation de d\'epart est $35 x + 38 y = 1$),
                        l'ensemble des solutions est
                              $\left\lbrace (38\,k-13 ; 12-35\,k) \mid k \in \mathbb{Z} \right\rbrace$.
            \end{ex}
            %
            \begin{thm}[Th\'eor\`eme de B\'ezout]
\index{Bezout@B\'ezout!Theoreme de@Th\'eor\`eme de}
                  Deux entiers relatifs $a$ et $b$ sont premiers entre eux si et seulement si
                  il existe deux entiers relatifs $u$ et $v$ tels que \thmbox{$a u + b v = 1$}.
            \end{thm}
            %
            \begin{prop}[Cons\'equences du lemme de Gauss]
                  \begin{enumerate}
                        \item Si un nombre premier $p$ divise un produit $ab$, alors $p$ divise $a$ ou $p$ divise $b$.
                        \item Si $a | n$, $b | n$ et $\pgcd{a}{b} = 1$, alors $ab | n$.
                  \end{enumerate}
            \end{prop}
            %
            \subsection*{\'Equations du type $a x + b y = c$}
            Soient $a$ et $b$ deux entiers relatifs non nuls, $d = \pgcd{a}{b}$ et $c \in \mathbb{Z}$.
            \begin{itemize}
                  \item L'\'equation $a x + b y = c$ a des solutions dans $\mathbb{Z}^2$ si et seulement si $d$ divise $c$.
                  \item Si $d | c$, on divise par $d$ : $a' x + b' y = c'$ avec $a = d a'$, $b = d b'$, $c = d c'$ et $\pgcd{a'}{b'} = 1$.
                        On cherche une solution particuli\`ere $(x_0 ; y_0)$ (par exemple en multipliant par $c'$ des coefficients de B\'ezout de $a'$ et $b'$),
                        puis le lemme de Gauss donne toutes les solutions :
                        \[
                              x = x_0 + b' k , \qquad y = y_0 - a' k , \qquad k \in \mathbb{Z} .
                        \]
            \end{itemize}
            %
            \begin{prop}[Inverses dans $\mathbb{Z}/n\mathbb{Z}$]
                  Soient $n \geq 2$ et $a\in\mathbb{Z}$.
                  \begin{enumerate}
                        \item Il existe $\bar{u}$ tel que $\bar{a} \cdot \bar{u} = \bar{1}$ dans $\mathbb{Z}/n\mathbb{Z}$
                              si et seulement si $\pgcd{a}{n} = 1$ (c'est le th\'eor\`eme de B\'ezout : $au + nv = 1$).
                              On peut alors r\'esoudre $\bar{a}\,x = \bar{b}$ en multipliant par $\bar{u}$ : $x = \bar{u}\,\bar{b}$.
                        \item Si $p$ est premier, tout \'el\'ement non nul de $\mathbb{Z}/p\mathbb{Z}$ a un inverse, et
                              $\bar{x} \cdot \bar{y} = \bar{0} \Rightarrow \left( \bar{x} = \bar{0} \text{ ou } \bar{y} = \bar{0} \right)$.
                              Ce n'est plus vrai si $n$ n'est pas premier : dans $\mathbb{Z}/6\mathbb{Z}$, $\bar{2}\cdot\bar{3} = \bar{0}$.
                  \end{enumerate}
            \end{prop}


\TCDLessonStyleOff
\clearpage
\providecommand{\congru}[3]{#1 \equiv #2 \quad \mathrm{mod}\ #3}
\providecommand{\base}[2]{\overline{#1}^{_{(#2)}}} % Fanoratana nombre #1 en base #2. En maths mode vao mandeha.
%%%%%%%%%%%%%%%%%%%%%%%%%%%%%%%%%%%%%%%%%%%%%%%%%%%%%%%%%%%%%%%%%%%%%%%%%%%%%%%%%%%%%%%%%%%%%%%%%%%%%%
\section{Exercices}
%
\begin{exo}\label{exo:arith-1}
	Pour chacun des couples $(a,b)$ suivants, effectuer la division euclidienne de $a$ par $b$ :
		\begin{enumerate}
			\item $(a,b) = (22,4)$,
			\item $(a,b) = (37,-5)$,
			\item $(a,b) = (-28,3)$,
			\item $(a,b) = (-93,-8)$,
			\item $(a,b) = (122,13)$,
			\item $(a,b) = (72,-11)$,
			\item $(a,b) = (-30,2)$,
			\item $(a,b) = (-45,-7)$.
		\end{enumerate}
\end{exo}
%
\begin{exo}\label{exo:arith-2}
	\begin{enumerate}
		\item Convertir le nombre $\base{100}{2}$ en base $ 10 $.
		\item Convertir le nombre $\base{1011}{2}$ en base $ 10 $.
		\item Convertir le nombre $\base{110010}{2}$ en base $ 10 $.
		\item Convertir le nombre $\base{1000101}{2}$ en base $ 10 $.
		\item Convertir le nombre $\base{12}{10}$ en base $ 2 $.
		\item Convertir le nombre $\base{237}{10}$ en base $ 2 $.
		\item Convertir le nombre $\base{154}{10}$ en base $ 2 $.
		\item Convertir le nombre $\base{73}{10}$ en base $ 2 $.
	\end{enumerate}
\end{exo}
%
\begin{exo}\label{exo:arith-3}
	\begin{enumerate}
		\item Convertir le nombre $\base{1222}{3}$ en base $ 10 $.
		\item Convertir le nombre $\base{1133}{5}$ en base $ 10 $.
		\item Convertir le nombre $\base{4035}{6}$ en base $ 10 $.
		\item Convertir le nombre $\base{11321}{8}$ en base $ 10 $.
		\item Convertir le nombre $\base{56}{10}$ en base $ 4 $.
		\item Convertir le nombre $\base{316}{10}$ en base $ 6 $.
		\item Convertir le nombre $\base{3317}{10}$ en base $ 7 $.
		\item Convertir le nombre $\base{9715}{10}$ en base $ 9 $.
	\end{enumerate}
\end{exo}
%
\begin{exo}\label{exo:arith-4}
	\begin{enumerate}
		\item Convertir le nombre $\base{122}{3}$ en base $ 5 $.
		\item Convertir le nombre $\base{31130}{4}$ en base $ 3 $.
		\item Convertir le nombre $\base{134334}{5}$ en base $ 7 $.
		\item Convertir le nombre $\base{203102}{6}$ en base $ 9 $.
		\item Convertir le nombre $\base{35006}{7}$ en base $ 6 $.
		\item Convertir le nombre $\base{7215}{8}$ en base $ 4 $.
	\end{enumerate}
\end{exo}
%
\begin{exo}\label{exo:arith-5}
      \begin{enumerate}
            \item D\'emontrer que la somme de trois entiers relatifs cons\'ecutifs est divisible par $3$.
            \item D\'emontrer que le produit de trois entiers relatifs cons\'ecutifs est divisible par $3$.
      \end{enumerate}
\end{exo}
%
\begin{exo}\label{exo:arith-6}
      Montrer que tout nombre premier sup\'erieur ou \'egal \`a $5$ est 
      soit de la forme $6 k + 1$ soit de la forme $6 k -1$ o\`u $k\in\mathbb{N}^*$.
\end{exo}
%
\begin{exo}\label{exo:arith-7}
      \begin{enumerate}
            \item D\'eterminer les entiers naturels non nuls $n$ dont 
                  la division par 16 donne un reste \'egal au carr\'e du quotient.
            \item D\'eterminer les entiers naturels non nuls $n$ dont
                  la division par 125 donne un reste \'egal au cube du quotient.
      \end{enumerate}
\end{exo}
%
\begin{exo}\label{exo:arith-8}
      \begin{enumerate}
            \item Soit $n$ un entier naturel sup\'erieur ou \'egal \`a 3. 
                  D\'eterminer le reste de la division de $n^3 - 2 $ par $n-2$.
            \item D\'eterminer les entiers naturels $n$ sup\'erieurs ou \'egal \`a 3 tel que $n -2 $ divise $n^3 - 2$.
      \end{enumerate}
\end{exo}
%
\begin{exo}\label{exo:arith-9}
      Soit $n$ un entier naturel et soit $r$ le reste de la division de $2^n $ par 7.
      \begin{enumerate}
            \item D\'eterminer le reste de la division de $2^{n+3}$ par 7 et en d\'eduire en fonction de $n$, 
                  le reste de la division de $2^n$ par 7.
            \item D\'eterminer le reste de la division de $23^{41}$ par 7.
      \end{enumerate}
\end{exo}
%
\begin{exo}\label{exo:arith-10}
      \begin{enumerate}
            \item D\'eterminer le reste de la division de $2^n$ par 7 pour $n \in \{ 1, \ 2, \ 3, \ 4, \ 5, \ 6 \}$.
            \item Montrer que 7 divise $2^{n+3} - 2^n$. 
                  En d\'eduire les valeurs de $n$ pour lesquelles la division de $2^n$ par 7 donne un reste \'egal \`a 1.
            \item Montrer que $9^n$ et $2^n$ ont m\^eme reste par la division par 7.
            \item Quel est le reste de la division de $2^{999} \cdot 9^{222}$ par 7\,?
      \end{enumerate}
\end{exo}
%
\begin{exo}\label{exo:arith-11}
      Soit $E$ l'ensemble des nombres entiers naturels $N$ qui s'\'ecrivent 
      sous la forme $N = \overline{ababab}$ en base 10, o\`u $a$ et $b$ sont des chiffres et $a$ est non nul.
      \begin{enumerate}
            \item Combien y a-t-il de nombres $N$ dans $E$\,? 
            \item Montrer que tous les nombres de $E$ sont multiples de $10101$.
      \end{enumerate}
\end{exo}
%
\begin{exo}\label{exo:arith-12}
      \begin{enumerate}
            \item Calculer pour $p = 0$, puis pour $p$ entier strictement positif, le reste de la division de $10^p$ par 6.
            \item Soit $n$ un entier sup\'erieur ou \'egal \`a 1. 
                  Un entier $x$ s'\'ecrit $\overline{a_n \, a_{n-1} \, \dots \, a_1 \, a_0} $ 
                  dans le syst\`eme de num\'eration en base 10.
                  D\'emontrer que 6 divise $x$ si et seulement si 6 divise $4 \; (a_n + a_{n-1} + \cdots + a_2 + a_1 ) + a_0$.
      \end{enumerate}
\end{exo}
%
\begin{exo}\label{exo:arith-13}
      \begin{enumerate}
            \item Soit $N$ un entier relatif impair. 
                  En \'ecrivant $N = 2\ k + 1 \ (k \in \mathbb{Z})$, d\'eterminer le reste de la division de $N^2$ par 8.
            \item R\'esoudre dans $\mathbb{Z} \times \mathbb{Z}$ 
                  l'\'equation $x^2 = 8 \, y + 1$ o\`u $x$ et $y$ sont les inconnues. 
            \item En d\'eduire que la parabole $(P)$ d'\'equation $y = \frac{x^2 - 1}{8}$ 
                  passe par une infinit\'e de points de coordonn\'ees enti\`eres.
      \end{enumerate}
\end{exo}
%
\begin{exo}\label{exo:arith-14}
      Soit $\{0,\ 1,\ 2,\ 3,\ 4,\ 5,\ 6,\ 7,\ 8,\ 9,\ A,\ B \}$ 
      l'ensemble des << chiffres >> qu'on utilise dans l'\'ecriture en base 12 des nombres ($A$ vaut dix et $B$ vaut onze).
      Un nombre $N$ s'\'ecrit $N = \overline{a_n \, a_{n-1} \, \dots \, a_1 \, a_0}^{_{(12)}}$ en base 12 $(n \geq 1)$.
      \begin{enumerate}
            \item Montrer que $\congru{N}{a_n + a_{n-1} + \cdots + a_2 + a_1  + a_0}{11}$.
            \item En d\'eduire un crit\`ere de divisibilit\'e par 11 d'un nombre $N$ \'ecrit en base 12.
            \item Un nombre $N$ s'\'ecrit $\overline{xy4}^{_{(12)}}$ en base 12 et $\overline{3yx}$ en base 10,
                  o\`u $x$ et $y$ sont des chiffres de $0$ \`a $9$, $x \neq 0$.
                  De plus, $N$ est divisible par 11.
                  D\'eterminer $N$.
      \end{enumerate}
\end{exo}
%
\begin{exo}\label{exo:arith-15}
      \begin{enumerate}
            \item V\'erifier que 7 divise les nombres $2^6 -1$ , $3^6 - 1$, $4^6 - 1$ et $5^6 - 1$.
            \item Pour tout entier naturel $n$, on d\'efinit le nombre $A_n$ par $A_n = 2^n + 3^n + 4^n + 5^n$.
                  Montrer que $A_{n+6} - A_n$ est divisible par 7.
            \item Soit $n$ un entier naturel, soient $q$ et $r$ son quotient et son reste respectivement par la division par 6.
                  Montrer que $A_n$ et $A_r$ ont m\^eme reste par la division par 7.
            \item D\'eterminer les valeurs de $n$ pour lesquelles $A_n$ est divisible par 7.
      \end{enumerate}
\end{exo}
%%%%%%%%%%%%%%%%%%%%%%%%%%%%%%%%%%%%%%%%%%%%%%%%%%%%%%%%%%%%%%%%%%%%%%%%%%%%%%%%%%%%
\begin{exo}\label{exo:arith-16}
      Trouver le PGCD de chacun des couples $(a;b)$ suivants :
      \begin{enumerate}
            \item $(a;b) = (499877;797111)$,
            \item $(a;b) = (97771;87567)$,
            \item $(a;b) = (27122;10097)$,
            \item $(a;b) = (15155;12958) $,
            \item $(a;b) = (15454;11345) $.
      \end{enumerate}
\end{exo}
%
\begin{exo}\label{exo:arith-17}
      Pour chacun des couples $(a;b)$ suivants, d\'eterminer un couple $(u;v)$ tel que
      $a\,u + b\,v = \pgcd{a}{b}$.
      \begin{enumerate}
            \item $(a;b) = (699 ; 624)$,
            \item $(a;b) = (972 ; 962)$,
            \item $(a;b) = (875 ; 788)$,
            \item $(a;b) = (934 ; 414)$,
            \item $(a;b) = (872 ; 432)$.
      \end{enumerate}
\end{exo}
%
\begin{exo}\label{exo:arith-18}
      R\'esoudre dans $\mathbb{Z}^2$ les \'equations suivantes : 
      \begin{enumerate}
            \item $3 u + 5 v = 1$,
            \item $13 u + 21 v = 1$,
            \item $37 u + 45 v = 1$,
            \item $3 u + 5 v = 2$,
            \item $13 u + 21 v = 3$,
            \item $37 u + 45 v = 4$,
            \item $6 u + 10 v = 6$,
            \item $39 u + 63 v = 6$.
      \end{enumerate}
\end{exo}
%
\begin{exo}\label{exo:arith-19}
      Soient $a$ et $b$ deux entiers relatifs non nuls et $n$ un entier naturel.
      \begin{enumerate}
            \item Montrer que si $a$ et $b$ sont premiers entre eux alors $5a + 3b$ et $3a + 2b$ sont aussi premiers entre eux.
            \item \begin{enumerate}
                  \item Montrer que $\pgcd{n^2 + 1}{n+1}$ divise $2$.
                  \item En d\'eduire que $n^2 + 1$ et $n+1$ sont premiers entre eux lorsque $n$ est pair.
                        Que vaut ce PGCD lorsque $n$ est impair\,?
                  \end{enumerate}
      \end{enumerate}
\end{exo}
%%%%%%%%%%%%%%%%%%%%%%%%%%%%%%%%%%%%%%%%%%%%%%%%%%%%%%%%%%%%%%%%%%%%%%%%%%%%%%%%%%%%
\begin{exo}\label{exo:arith-20}% ~ \newline
\begin{enumerate}
\item D\'eterminer l'ensemble des entiers relatifs $x$ tels que : $123 \cdot x \equiv 456 \quad [7]$.
\item \begin{enumerate}
      \item D\'ecomposer 78 et 14 en produit de facteurs premiers.
      \item En d\'eduire dans $\mathbb{N}$ l'ensemble des diviseurs de 78 et l'ensemble des diviseurs de 14.
      \item D\'eterminer alors $g = PGCD (78 , 14)$ et $m = PPCM (78 , 14)$.
      \end{enumerate}
\end{enumerate}
\end{exo}
%%%%%%%%%%%%%%%%%%%%%%%%%%%%%%%%%%%%%%%%%%%%%%%%%%%%%%%%%%%%%%%%%%%%%%%%%%%%%%%%%%%%
\begin{exo}\label{exo:arith-21}%~\newline
\begin{enumerate}
\item Trouver tous les nombres entiers naturels ayant 3 chiffres et divisibles \`a la fois par 14 et par 34.
\item Trouver deux entiers naturels $a$ et $b$ tels que : 
      $$
      \left\lbrace
      \begin{array}{l}
      PGCD (a,b) = 6 \\
      a \cdot b = 2700      .
      \end{array}
      \right.
      \;
      $$
      En d\'eduire le $PPCM$ de $a$ et $b$.
\end{enumerate}
\end{exo}
%%%%%%%%%%%%%%%%%%%%%%%%%%%%%%%%%%%%%%%%%%%%%%%%%%%%%%%%%%%%%%%%%%%%%%%%%%%%%%%%%%%%
\begin{exo}\label{exo:arith-22}
      \begin{enumerate}
            \item Montrer par r\'ecurrence que $9^n - 2^n$ est divisible par $7$ pour tout $n \in \mathbb{N}$.
            \item Montrer que, pour tout entier naturel $n$ non-nul, $3n-2$ et $5n-3$ sont premiers entre eux.
            \item On consid\`ere dans $\mathbb{Z} / 7 \mathbb{Z}$, l'\'equation : 
                  \begin{equation}
                  \TCDtag{E}
                        x^2 - x + 1  = 0 .
                  \label{eq:E arithm}
                  \end{equation}
                  \begin{enumerate}
                        \item V\'erifier que $5$ est une solution de \eqref{eq:E arithm}.
                        \item D\'emontrer que \eqref{eq:E arithm} est \'equivalente \`a : 
                              $(x-4)^2 = 1$.
                        \item En d\'eduire l'ensemble $S$ des solutions de \eqref{eq:E arithm}.
                  \end{enumerate}
      \end{enumerate}
\end{exo}
%
\begin{exo}\label{exo:arith-23}
      \begin{enumerate}
            \item Calculer $PGCD(319,232,145)$ puis r\'esoudre dans $\mathbb{Z} \times \mathbb{Z}$
                  l'\'equation $319 x + 232 y = 145$.
            \item R\'esoudre l'\'equation $13x - 7 y = 11$,
                  d'inconnues $(x,y) \in \mathbb{N} \times \mathbb{N}$.
      \end{enumerate}
\end{exo}

\begin{exo}\label{exo:arith-24}
	Montrer que, pour tout entier relatif $n$, $n^2 + n + 2$ n'est pas divisible par $3$.
\end{exo}
%
\begin{exo}\label{exo:arith-25}
      \begin{enumerate}
            \item \begin{enumerate}
                        \item \'Ecrire en base $10$ le nombre $N = \base{21012}{3}$.
                        \item \'Ecrire le nombre $100$ en base $3$.
                        \item Comment reconna\^it-on, sur son \'ecriture en base $3$, qu'un entier naturel est divisible par $3$\,?
                              par $9$\,?
                  \end{enumerate}
            \item \begin{enumerate}
                        \item Dresser les tables d'addition et de multiplication de $\mathbb{Z}/3\mathbb{Z}$.
                        \item Montrer que l'\'equation $x^2 + x + \bar{2} = \bar{0}$ n'a pas de solution dans $\mathbb{Z}/3\mathbb{Z}$.
                        \item R\'esoudre dans $\mathbb{Z}/3\mathbb{Z}$ l'\'equation $x^2 + \bar{2}\,x + \bar{1} = \bar{0}$.
                        \item R\'esoudre dans $\mathbb{Z}/3\mathbb{Z} \times \mathbb{Z}/3\mathbb{Z}$ le syst\`eme
                              \[
                                    \left\lbrace
                                          \begin{array}{l}
                                                x + y = \bar{2} \\
                                                x + \bar{2}\, y = \bar{0} \;.
                                          \end{array}
                                    \right.
                              \]
                  \end{enumerate}
      \end{enumerate}
\end{exo}
%
\begin{exo}\label{exo:arith-26}
      \begin{enumerate}
            \item Dresser la table des carr\'es des \'el\'ements de $\mathbb{Z}/5\mathbb{Z}$.
            \item R\'esoudre dans $\mathbb{Z}/5\mathbb{Z}$ les \'equations $x^2 = \bar{4}$ et $x^2 = \bar{2}$.
                  En d\'eduire que l'\'ecriture d\'ecimale du carr\'e d'un entier ne se termine jamais par $2$, $3$, $7$ ou $8$.
            \item Calculer $x^4$ pour tout $x \in \mathbb{Z}/5\mathbb{Z}$.
                  En d\'eduire que $x^5 = x$ pour tout $x \in \mathbb{Z}/5\mathbb{Z}$.
            \item En d\'eduire que, pour tout entier naturel $n$, $n^5 - n$ est divisible par $5$.
            \item \begin{enumerate}
                        \item V\'erifier que $n^5 - n = (n-1)\, n\, (n+1)\, (n^2 + 1)$.
                        \item Montrer que $n^5 - n$ est divisible par $6$, puis par $30$.
                  \end{enumerate}
            \item Application : montrer que, pour tout entier naturel $n$, les nombres $n$ et $n^5$ ont le m\^eme chiffre des unit\'es.
      \end{enumerate}
\end{exo}
%%%%%%%%%%%%%%%%%%%%%%%%%%%%%%%%%%%%%%%%%%%%%%%%%%%%%%%%%%%%%%%%%%%%%
%%%%%%%%%%%%%%%%%%%%%%%%%%%%%%%%%%%%%%%%%%%%%%%%%%%%%%%%%%%%%%%%%%%%%
\begin{exo}\label{exo:arith-27}
	On dispose de deux r\'ecipients : l'un peut mesurer 5 litres tandis que l'autre peut en mesurer 3.
	Peut-on, \`a l'aide de ces deux r\'ecipients, mesurer 1 litre\,? Si oui, comment\,?
\end{exo}
%
\begin{exo}\label{exo:arith-28}
	Andry poss\`ede un tas de pi\`eces de 4~Ar et Bema poss\`ede un tas de pi\`eces de 10~Ar.
	\begin{enumerate}
		\item Serait-il possible pour Andry de donner 18~Ar \`a Bema\,?
		\item Serait-il possible pour Bema de donner 25~Ar \`a Andry\,?
	\end{enumerate}
\end{exo}
%
\begin{exo}\label{exo:arith-29}
	Sauriez-vous me dire quel est le dernier chiffre de chacun des nombres suivants\,?
		\[ 
			10^{10}			,\qquad
			5^{5^{5}}			,\qquad
			2^{3^{4^{5}}} 	.
		\]
\end{exo}
%
\begin{exo}\label{exo:arith-30}
	Quel est le plus grand diviseur commun de chacun des couples de nombres suivants :
	\[
		(4,12)				,\qquad
		(60,24)			,\qquad
		(9379,8039)		.
	\]
\end{exo}
%
\begin{exo}\label{exo:arith-31}
	\begin{enumerate}
		\item 
				Jean le Baptiste partit dans le d\'esert un dimanche.
			Il y resta 40 jours.
			Quel jour de la semaine Jean est-il revenu du d\'esert\,?
			
				Quel jour de la semaine Jean serait-il revenu s'il y \'etait rest\'e 40~000 jours\,? 40~000~000 de jours\,?
		\item 
			Un dimanche, un visiteur frappa \`a ma porte.
				Il me raconta qu'il revenait d'un voyage de 50 jours.
				Quel jour de la semaine avait-il commenc\'e son voyage\,?
			
				Quel jour aurait-ce \'et\'e s'il revenait d'un voyage de 50~000 jours\,? De 50~000~000 de jours\,?
	\end{enumerate}
\end{exo}
%
\begin{exo}\label{exo:arith-32}
		Rakoto se trouve incapable de payer son loyer du mois de janvier.
		Il poss\`ede cependant une barre d'argent de 31~cm de long.
		Il alla parler au propri\'etaire et ils trouv\`erent un arrangement :
		\`a la fin de chaque jour du mois de janvier, le propri\'etaire devra d\'etenir, comme gage,
		une longueur totale d'argent \'egale au num\'ero du jour.
	De cette mani\`ere, au $n$-ième jour du mois de janvier, le propri\'etaire aura re\c{c}u $n$~cm d'argent.
		Ils continueront ainsi tout au long du mois, jusqu'\`a ce que Rakoto trouve un moyen de payer son loyer.
	
		Il se trouve que Rakoto ne veut pas d\'ecouper sa barre en trop de morceaux ;
	il cherche \`a la d\'ecouper en le plus petit nombre possible de morceaux qui lui permette de tenir son engagement
	(chaque jour, il peut reprendre des morceaux d\'ej\`a donn\'es et en donner d'autres).
		Comment Rakoto doit-il diviser sa barre d'argent\,?
%	
%	\noindent\textsc{Exemple :}\newline
%		Il peut le diviser en morceaux de longueurs respectifs : 1~cm, 2~cm, 3~cm, 4~cm, 4~cm, 5~cm, 6~cm et 7~cm ;
%		soit 8 morceaux.
%		\begin{enumerate}
%			\item Au premier jour, il donnera le morceau de 1~cm,
%			\item le second jour, il reprendra le morceau de 1~cm et donnera le morceau de 2~cm,
%			\item il continuera de m\^eme jusqu'au setpi\`eme jours.
%			\item Au huiti\`eme jour, il ajoutera le morceau de 1~cm au morceau de 7~cm qu'il ne reprendra pas,
%			\item et ainsi de suite.
%		\end{enumerate}
\end{exo}

\begin{exo}\label{exo:arith-33}
      Pour tout entier naturel $n$, on consid\`ere la fraction $\dfrac{n+7}{2n+5}$.
      \begin{enumerate}
            \item Montrer que tout diviseur commun de $n+7$ et $2n+5$ divise $9$.
            \item Montrer que $3$ divise $n + 7$ si et seulement si $n \equiv 2 \ [3]$, et que dans ce cas $3$ divise aussi $2n+5$.
            \item En d\'eduire les entiers naturels $n$ pour lesquels la fraction est irr\'eductible.
            \item Pour quels entiers naturels $n$ peut-on simplifier la fraction par $9$\,?
                  Simplifier la fraction pour $n = 20$.
      \end{enumerate}
\end{exo}

\begin{exo}\label{exo:arith-34}
      On cherche le plus petit entier naturel $N$ dont l'\'ecriture d\'ecimale commence par le chiffre $1$
      et tel que, si l'on d\'eplace ce chiffre $1$ \`a la fin de l'\'ecriture, on obtient le triple de $N$.
      Par exemple, si $N$ \'etait $1\,234$, le nombre obtenu serait $2\,341$.
      \begin{enumerate}
            \item On note $k$ le nombre de chiffres qui suivent le $1$, de sorte que $N = 10^k + x$ avec $0 \leq x < 10^k$.
                  Exprimer le nombre obtenu en fonction de $x$ et montrer que la condition s'\'ecrit $7 x = 3 \cdot 10^k - 1$.
            \item D\'eterminer le reste de la division de $10^k$ par $7$ pour $k \in \{1, 2, \ldots, 6\}$.
            \item Montrer que $7$ divise $3 \cdot 10^k - 1$ si et seulement si $10^k \equiv 5 \ [7]$.
            \item En d\'eduire le plus petit nombre $N$ cherch\'e et v\'erifier le r\'esultat.
      \end{enumerate}
\end{exo}

\begin{exo}\label{exo:arith-35}
      \begin{enumerate}
            \item D\'eterminer les restes de la division de $5^n$ par $13$ pour $n \in \{0, 1, 2, 3, 4\}$.
                  En d\'eduire, selon les valeurs de $n$, le reste de la division de $5^n$ par $13$.
            \item D\'eterminer les entiers naturels $n$ tels que $13$ divise $5^n - 1$, puis ceux tels que $13$ divise $5^n + 1$.
            \item Montrer que $5^n + 2$ n'est divisible par $13$ pour aucun entier naturel $n$.
            \item Quel est le reste de la division de $5^{2026} + 5^{2027}$ par $13$\,?
      \end{enumerate}
\end{exo}

\begin{exo}\label{exo:arith-36}
      On consid\`ere les suites $(u_n)$ et $(v_n)$ d\'efinies pour tout entier naturel $n$ par
      $u_n = 3 + 5n$ et $v_n = 7 + 8n$.
      On note $A = \{u_0, u_1, \ldots, u_{99}\}$ et $B = \{v_0, v_1, \ldots, v_{99}\}$.
      \begin{enumerate}
            \item Montrer qu'un nombre appartient \`a la fois \`a $A$ et \`a $B$ si et seulement si il s'\'ecrit
                  $3 + 5p = 7 + 8q$ avec $p$ et $q$ entiers compris entre $0$ et $99$.
            \item R\'esoudre dans $\mathbb{Z}^2$ l'\'equation $5p - 8q = 4$.
            \item En d\'eduire que les nombres communs \`a $A$ et $B$ sont les termes d'une suite arithm\'etique
                  dont on donnera le premier terme et la raison. Comparer la raison \`a $PPCM(5;8)$.
            \item Combien de nombres appartiennent \`a $A \cap B$\,?
                  Combien de nombres distincts y a-t-il dans $A \cup B$\,?
      \end{enumerate}
\end{exo}

\begin{exo}\label{exo:arith-37}
      On cherche les triplets $(a, b, c)$ d'entiers naturels tels que
      $a < b < c$, $a$ divise $b$, $b$ divise $c$ et $a + b + c = 55$.
      \begin{enumerate}
            \item Montrer que $a$ divise $55$.
            \item On \'ecrit $b = k a$ et $c = m b$. Montrer que $k \geq 2$, $m \geq 2$ et $a\,(1 + k + k m) = 55$.
            \item D\'eterminer tous les triplets cherch\'es.
            \item Quelle est la plus grande valeur possible de $c$\,?
      \end{enumerate}
\end{exo}

\begin{exo}\label{exo:arith-38}
      \begin{enumerate}
            \item V\'erifier que, pour tout entier relatif $n$, on a $n^2 + 7 = (n+3)(n-3) + 16$.
            \item En d\'eduire que, pour $n \neq -3$, l'entier $n+3$ divise $n^2 + 7$ si et seulement si $n + 3$ divise $16$.
            \item D\'eterminer tous les entiers relatifs $n$ tels que $n + 3$ divise $n^2 + 7$.
            \item Quel est le plus grand d'entre eux\,? Donner alors le quotient de $n^2 + 7$ par $n + 3$.
      \end{enumerate}
\end{exo}

\begin{exo}\label{exo:arith-39}
      Pour un entier $n \geq 3$, on s'int\'eresse aux fractions irr\'eductibles parmi
      $\dfrac{1}{n},\ \dfrac{2}{n},\ \ldots ,\ \dfrac{n-1}{n}$.
      \begin{enumerate}
            \item Pour $n = 12$, \'ecrire la liste de ces fractions irr\'eductibles et calculer leur somme.
            \item Montrer que, pour $1 \leq k \leq n-1$, on a $\pgcd{k}{n} = \pgcd{n-k}{n}$.
            \item En d\'eduire que si $\dfrac{k}{n}$ est irr\'eductible, alors $\dfrac{n-k}{n}$ l'est aussi,
                  et que ces deux fractions sont diff\'erentes.
            \item On note $N$ le nombre de ces fractions irr\'eductibles et $S$ leur somme.
                  En regroupant les fractions deux par deux, montrer que $S = \dfrac{N}{2}$.
                  En d\'eduire que $N$ est pair.
            \item Application : calculer $N$ et $S$ pour $n = 15$.
      \end{enumerate}
\end{exo}

\begin{exo}[Nombres de Mersenne]\label{exo:arith-40}
	On appelle \emph{nombres de Mersenne} les nombres de la forme $2^n - 1$ o\`u $n \in \mathbb{N}^*$.
	\begin{enumerate}
		\item D\'emontrer que $2^q - 1$ divise $2^{pq} - 1$ pour tout $(p,q) \in (\mathbb{N}^*)^2$.
		\item En d\'eduire que si $2^n - 1$ est premier alors $n$ est premier.
		\item Trouver un nombre premier $p$ tel que $2^p - 1$ n'est pas premier.
	\end{enumerate}
\end{exo}
% IO AMBONY IHANY ITO
%\begin{exo}
%	Pour chaucune des propositions suivantes, dire si elle est vraie ou fausse ; justifier votre r\'eponse dans chaque cas.
%	\begin{itemize}
%		\item Si $n$ est premier alors nombre $2^n - 1$ est premier.
%		\item Si $n$ est n'est pas premier alors nombre $2^n - 1$ n'est pas premier.
%	\end{itemize}
%	Note : Le nombre $n$ consid\'er\'e est suppos\'e entier naturel non-nul.
%\end{exo}
%
\begin{exo}[Propri\'et\'es des nombres de Fermat]\label{exo:arith-42}
\index{Fermat!nombre de}
	On d\'efinit le \emph{nombre de Fermat} de rang $n$ par $F_n := 2^{2^n} + 1$ o\`u $n \in \mathbb{N}$.
	\begin{enumerate}
		\item
			\begin{enumerate}
				\item D\'emontrer que $\prod_{k = 0}^{n -1} F_k := F_0 \cdot F_1 \cdots F_{n-1}$ 
					est \'egal \`a $F_n - 2$ pour tout $n\in\mathbb{N}^*$.
				\item En d\'eduire que les nombres de Fermat sont deux \`a deux premiers entre eux.
			\end{enumerate}
		\item Montrer que $F_n$ divise $2^{F_n} - 2$ pour tout $n$.
	\end{enumerate}
%	
%	Pierre de Fermat (1601 -- 1665) \'etait un magistrat fran\c{c}ais et math\'ematicien amateur.
%	Fermat, apr\`es avoir v\'erifi\'e que $F_i$ \'etait premier pour $i \in \{ 0,\ 1,\ 2,\ 3 ,\ 4 \}$,
%	conjecturat que $F_n$ \'etait premier pour tout $n$.
%	Leonhard Euler (1707 -- 1783), l'illustre math\'ematicien suisse, d\'emontrat que $F_5$ est divisible par $641$, 
%	et par cons\'equent n'est pas premier.
\end{exo}

\begin{exo}\label{exo:arith-43}
	\begin{enumerate}
		\item 
			\begin{enumerate}
				\item Montrer que pour tout nombre pair $a \geq 0$, $3^a \equiv 1 \ [4]$.
					\label{quest: Arithm Exo a}
					\item Montrer que $2^{4k} \equiv 1 \ [5]$, pour tout $k \in \mathbb{N}$.
					\label{quest: Arithm Exo b}
				\item D\'eduire des questions \ref{quest: Arithm Exo a} et \ref{quest: Arithm Exo b} que 
					pour tout nombre pair $a \geq 0$, $2^{3^a} \equiv 2 \ [10]$.
				\item Utiliser ce dernier r\'esultat pour donner le dernier chiffre de l'\'ecriture d\'ecimale du nombre
					$2^{3^{4^5}}$.
			\end{enumerate}
		\item Donner les deux derniers chiffres de l'\'ecriture d\'ecimale du nombre $7^{7^{7^{7^{7^{7^7}}}}}$.
		\item Quels sont les trois derniers chiffres de l'\'ecriture d\'ecimale du nombre : 
			$\displaystyle 2^{3^{4^{5}}} \times 5^{4^{3^{2}}}$. \trollface
	\end{enumerate}
\end{exo}

\begin{exo}\label{exo:arith-44}
	Montrer qu'il existe (au moins) un terme de la suite :
		\[
			3			\, ;\ 
			33			\, ;\ 
			333			\, ;\ 
			3\,333		\, ;\ 
			\ldots
		\]
		qui est divisible par $2\,019$.
\end{exo}

\begin{exo}[Ne pensez pas trop loin \trollface]\label{exo:arith-45}
	Donner \emph{tous} les couples de nombres premiers $( x,y )$, $x < y$, tels que $x + y = 2\,019$.
\end{exo}

\begin{exo}\label{exo:arith-46}
	\begin{enumerate}
		\item Soit $P(n) = n^2 + n + 41$.
			Donner cinq entiers naturels $n_1$, $n_2$, $n_3$, $n_4$ et $n_5$ tels 
			qu'aucun des nombres $P(n_1)$, $P(n_2)$, $P(n_3)$, $P(n_4)$ et $P(n_5)$ n'est premier.
			\item D\'emontrer que si $p$ est un nombre premier sup\'erieur ou \'egal \`a $5$ alors $p^2 \equiv 1 \quad [24]$.
	\end{enumerate}
\end{exo}
























\begin{corrige}
\section{Solutions des exercices et preuves des th\'eor\`emes}
\begin{proof}
(Preuve de la Propri\'et\'e~\ref{prop:divisibilite}.)
Dans toute la preuve, $a$, $b$ et $c$ sont des entiers relatifs non nuls.
\begin{enumerate}
\item Supposons que $a|b$ et $b|a$.
      Il existe $q \in \mathbb{Z}$ tel que $b = q a$ et $q' \in \mathbb{Z}$ tel que $a = q' b$.
      En combinant ces deux \'egalit\'es, on obtient : $b = q ( q' b) = (q q') b$, soit $(qq' - 1)\, b = 0$.
      Comme $b \neq 0$ (hypoth\`ese de la propri\'et\'e), il suit que $qq' = 1$.
      (Si l'on autorisait $b = 0$, on aurait $a = q' b = 0$, donc encore $a = b$.)
      Comme $q$ et $q'$ sont des entiers dont le produit vaut $1$, il y a deux possibilit\'es :
            $q = q' = 1$ ou $q = q' = -1$.
      Dans le premier cas, on a $b = a$ ; dans le second cas, on a $b = - a$.
\item \begin{enumerate}
      \item $a = 1 \cdot a$, donc $a | a$.
      \item Supposons que $a | b $ et que $b | c$.
            Il existe deux entiers $k$ et $k'$ tels que $b = ak$ et $c = bk'$.
            Alors $c = (ak)k' = a(kk')$ avec $kk' \in \mathbb{Z}$, ce qui prouve que $a | c$.
      \item D'apr\`es le point 1, $a = b$ ou $a = -b$, donc $|a| = |b|$.
      \end{enumerate}
\item \begin{enumerate}
      \item $a | b$ signifie qu'il existe un entier $k$ tel que $b = ak$.
            En multipliant cette \'egalit\'e par $c$, on a $b c = a (kc)$ avec $kc \in \mathbb{Z}$, 
            ce qui prouve que $a | bc$.
      \item Avec $b = ak$, on a $bc = (ac) k$, ce qui montre que $ac | bc$.
      \item Si $b = ak$ et $c = ak'$, alors pour tout $(\beta, \gamma) \in \mathbb{Z}^2$,
            $\beta b + \gamma c = a(\beta k + \gamma k')$, donc $a | (\beta b + \gamma c)$.
      \end{enumerate}
\item $a = a \cdot 1$, donc $1 | a$ ; $0 = a \cdot 0$, donc $a | 0$.
\item Si $b | a$, il existe $k \in \mathbb{Z}$ tel que $a = kb$. Comme $a \neq 0$, on a $k \neq 0$ et $b \neq 0$, donc $|k| \geq 1$ et
      $|a| = |k|\,|b| \geq |b| > 0$.
\end{enumerate}
\end{proof}
%
\subsection*{Exercices de la le\c{c}on}
%
\begin{sol}[Exercice~\ref{exe:arith-1}]
On teste les nombres premiers inf\'erieurs ou \'egaux \`a la racine carr\'ee.
\begin{itemize}
\item $527 = 17 \times 31$ : $527$ n'est pas premier.
\item $\sqrt{769} \simeq 27{,}7$ ; aucun des nombres premiers $2, 3, 5, 7, 11, 13, 17, 19, 23$ ne divise $769$ : $769$ est premier.
\item $\sqrt{1249} \simeq 35{,}3$ ; aucun des nombres premiers $2, 3, \ldots, 31$ ne divise $1249$ : $1249$ est premier.
\item $57 = 3 \times 19$ : $57$ n'est pas premier (la somme de ses chiffres, $12$, est un multiple de $3$).
\end{itemize}
\end{sol}
%
\begin{sol}[Exercice~\ref{exe:arith-2}]
$\overline{101110}^{_{(2)}} = 2^5 + 2^3 + 2^2 + 2 = 46$ ;
$\overline{756}^{_{(8)}} = 7 \cdot 64 + 5 \cdot 8 + 6 = 494$ ;
$\overline{B00A16}^{_{(12)}} = 11 \cdot 12^5 + 10 \cdot 12^2 + 1 \cdot 12 + 6 = 2\,737\,152 + 1\,440 + 12 + 6 = 2\,738\,610$.
\end{sol}
%
\begin{sol}[Exercice~\ref{exe:arith-3}]
\begin{enumerate}
\item $5^2 = 25 \equiv -1 \ [13]$, donc $5^4 \equiv 1 \ [13]$. Comme $2018 = 4 \times 504 + 2$,
      $5^{2018} = (5^4)^{504} \cdot 5^2 \equiv 1 \cdot 25 \equiv 12 \ [13]$. Le reste est $\boxed{12}$.
\item $2^3 = 8 \equiv 1 \ [7]$ et $2019 = 3 \times 673$, donc $2^{2019} = (2^3)^{673} \equiv 1 \ [7]$. Le reste est $\boxed{1}$.
\item $5^3 = 125 = 7 \times 17 + 6 \equiv -1 \ [7]$, donc $5^6 \equiv 1 \ [7]$.
      Comme $4037 = 6 \times 672 + 5$, $5^{4037} \equiv 5^5 = 5^3 \cdot 5^2 \equiv (-1) \cdot 25 \equiv -25 \equiv 3 \ [7]$.
      Le reste est $\boxed{3}$.
\end{enumerate}
\end{sol}
%
\begin{sol}[Exercice~\ref{exe:arith-4}]
On a $N = a_0 + 10\,(a_1 + 10\, a_2 + \cdots + 10^{n-1} a_n)$.
\begin{enumerate}
\item Comme $10 \equiv 0 \ [2]$, on a $N \equiv a_0 \ [2]$ ; donc $2 | N \iff a_0 \equiv 0 \ [2]$
      (le chiffre des unit\'es est $0$, $2$, $4$, $6$ ou $8$).
\item De m\^eme $10 \equiv 0 \ [5]$, donc $N \equiv a_0 \ [5]$ et $5 | N \iff a_0 \equiv 0 \ [5]$ ($a_0 = 0$ ou $5$).
\item $10 \equiv 1 \ [3]$, donc $10^p \equiv 1 \ [3]$ pour tout $p$, puis $a_p 10^p \equiv a_p \ [3]$ ;
      en sommant, $N \equiv \sum_{p=0}^n a_p \ [3]$, d'o\`u $3 | N \iff \sum_{p=0}^n a_p \equiv 0 \ [3]$.
\end{enumerate}
\end{sol}
%
\begin{sol}[Exercice~\ref{exe:arith-5}]
Dans $\mathbb{Z}/7\mathbb{Z}$ : $\bar{2}\cdot\bar{4} = \bar{1}$, $\bar{3}\cdot\bar{5} = \bar{1}$, $\bar{6}\cdot\bar{6} = \bar{1}$.
\begin{enumerate}
\item \begin{enumerate}
      \item $\bar{3}x = \bar{5} - \bar{4} = \bar{1}$ ; en multipliant par $\bar{5}$ : $x = \bar{5}$.
      \item $\bar{2}x = \bar{3} - \bar{6} = \bar{4}$ ; en multipliant par $\bar{4}$ : $x = \bar{16} = \bar{2}$.
      \item $\bar{6}x = \bar{1} - \bar{5} = \bar{3}$ ; en multipliant par $\bar{6}$ : $x = \bar{18} = \bar{4}$.
      \end{enumerate}
\item \begin{enumerate}
      \item En multipliant la premi\`ere \'equation par $\bar{4}$ : $x + \bar{6}y = \bar{5}$, soit $x = \bar{5} + y$.
            La seconde devient $\bar{3}(\bar{5} + y) + \bar{2}y = \bar{1}$, soit $\bar{1} + \bar{5}y = \bar{1}$, donc $\bar{5} y = \bar{0}$ et $y = \bar{0}$.
            Solution unique : $(x ; y) = (\bar{5} ; \bar{0})$.
      \item $x = \bar{6} - \bar{2}y$ ; la seconde donne $\bar{12} - \bar{4}y + \bar{3}y = \bar{4}$, soit $-y = \bar{4} - \bar{5} = -\bar{1}$,
            donc $y = \bar{1}$ et $x = \bar{4}$. Solution unique : $(\bar{4} ; \bar{1})$.
      \item $y = \bar{3} - \bar{5}x = \bar{3} + \bar{2}x$ ; la seconde donne $\bar{3}x + \bar{6} + \bar{4}x = \bar{0}$, soit
            $\bar{7}x + \bar{6} = \bar{6} = \bar{0}$ : c'est impossible. Le syst\`eme n'a \textbf{aucune} solution.
      \item $y = \bar{5} - \bar{3}x = \bar{5} + \bar{4}x$ ; la premi\`ere donne $\bar{2}x + \bar{15} + \bar{12}x = \bar{1}$, soit
            $\bar{14}x + \bar{1} = \bar{1}$, ce qui est toujours vrai. Les deux \'equations sont \'equivalentes
            (la premi\`ere est $\bar{3}$ fois la seconde) et l'ensemble des solutions est
            $\left\lbrace (x ; \bar{5} + \bar{4}x) \mid x \in \mathbb{Z}/7\mathbb{Z} \right\rbrace$, soit sept couples :
            $(\bar{0};\bar{5})$, $(\bar{1};\bar{2})$, $(\bar{2};\bar{6})$, $(\bar{3};\bar{3})$, $(\bar{4};\bar{0})$, $(\bar{5};\bar{4})$, $(\bar{6};\bar{1})$.
      \end{enumerate}
\end{enumerate}
\end{sol}
%
\begin{sol}[Exercice~\ref{exe:arith-6}]
\begin{enumerate}[label=\alph*)]
\item Suite des restes de l'algorithme d'Euclide : $4062$, $3215$, $847$, $674$, $173$, $155$, $18$, $11$, $7$, $4$, $3$, $1$, $0$.
      Donc $\pgcd{51\,959}{47\,897} = 1$.
\item $24\,881 = 1 \cdot 23\,213 + 1\,668$ ; $23\,213 = 13 \cdot 1\,668 + 1\,529$ ; $1\,668 = 1 \cdot 1\,529 + 139$ ; $1\,529 = 11 \cdot 139$.
      Donc $\pgcd{23\,213}{24\,881} = 139$.
\item $62\,651 = 1 \cdot 47\,027 + 15\,624$ ; $47\,027 = 3 \cdot 15\,624 + 155$ ; $15\,624 = 100 \cdot 155 + 124$ ;
      $155 = 1 \cdot 124 + 31$ ; $124 = 4 \cdot 31$. Donc $\pgcd{47\,027}{62\,651} = 31$.
\end{enumerate}
\end{sol}
%
\begin{sol}[Exercice~\ref{exe:arith-7}]
\begin{enumerate}
\item $615\,627 = 17 \cdot 36\,183 + 516$ ; $36\,183 = 70 \cdot 516 + 63$ ; $516 = 8 \cdot 63 + 12$ ; $63 = 5 \cdot 12 + 3$ ; $12 = 4 \cdot 3$.
      Donc $\pgcd{615\,627}{36\,183} = 3$. L'algorithme d'Euclide \'etendu donne
      $615\,627 \times (-2\,875) + 36\,183 \times 48\,916 = 3$.
\item $210\,573 = 23 \cdot 9\,099 + 1\,296$ ; $9\,099 = 7 \cdot 1\,296 + 27$ ; $1\,296 = 48 \cdot 27$. Donc $\pgcd{210\,573}{9\,099} = 27$.
      En remontant : $27 = 9\,099 - 7 \cdot 1\,296 = 9\,099 - 7\,(210\,573 - 23 \cdot 9\,099) = -7 \cdot 210\,573 + 162 \cdot 9\,099$.
      Donc $(u ; v) = (-7 ; 162)$.
\item Suite des restes : $3\,911$, $296$, $63$, $44$, $19$, $6$, $1$, $0$ ; donc $\pgcd{12\,325}{4\,207} = 1$.
\item Suite des restes (en divisant $23\,369$ par $13\,567$) : $9\,802$, $3\,765$, $2\,272$, $1\,493$, $779$, $714$, $65$, $64$, $1$, $0$ ;
      donc $\pgcd{13\,567}{23\,369} = 1$.
\end{enumerate}
Les coefficients de B\'ezout ne sont pas uniques : $(u + 12\,061\,k ; v - 205\,209\,k)$ convient aussi dans le premier cas.
\end{sol}
%
\subsection*{Exercices de fin de chapitre}
%
\begin{sol}[Exercice~\ref{exo:arith-1}]
On cherche $(q, r)$ avec $a = bq + r$ et $0 \leq r < |b|$.
\begin{center}
\begin{tabular}{|c|c|c|c|}
\hline $(a,b)$ & \'egalit\'e & $q$ & $r$ \\ \hline
$(22,4)$ & $22 = 4 \cdot 5 + 2$ & $5$ & $2$ \\
$(37,-5)$ & $37 = (-5)\cdot(-7) + 2$ & $-7$ & $2$ \\
$(-28,3)$ & $-28 = 3 \cdot (-10) + 2$ & $-10$ & $2$ \\
$(-93,-8)$ & $-93 = (-8) \cdot 12 + 3$ & $12$ & $3$ \\
$(122,13)$ & $122 = 13 \cdot 9 + 5$ & $9$ & $5$ \\
$(72,-11)$ & $72 = (-11)\cdot(-6) + 6$ & $-6$ & $6$ \\
$(-30,2)$ & $-30 = 2 \cdot (-15) + 0$ & $-15$ & $0$ \\
$(-45,-7)$ & $-45 = (-7) \cdot 7 + 4$ & $7$ & $4$ \\ \hline
\end{tabular}
\end{center}
\end{sol}
%
\begin{sol}[Exercice~\ref{exo:arith-2}]
$\base{100}{2} = 4$ ; $\base{1011}{2} = 11$ ; $\base{110010}{2} = 50$ ; $\base{1000101}{2} = 69$ ;
$12 = \base{1100}{2}$ ; $237 = \base{11101101}{2}$ ; $154 = \base{10011010}{2}$ ; $73 = \base{1001001}{2}$.
\end{sol}
%
\begin{sol}[Exercice~\ref{exo:arith-3}]
$\base{1222}{3} = 27 + 18 + 6 + 2 = 53$ ; $\base{1133}{5} = 125 + 25 + 15 + 3 = 168$ ;
$\base{4035}{6} = 864 + 18 + 5 = 887$ ; $\base{11321}{8} = 4096 + 512 + 192 + 16 + 1 = 4817$ ;
$56 = \base{320}{4}$ ; $316 = \base{1244}{6}$ ; $3317 = \base{12446}{7}$ ; $9715 = \base{14284}{9}$.
\end{sol}
%
\begin{sol}[Exercice~\ref{exo:arith-4}]
On passe par la base $10$ :
$\base{122}{3} = 17 = \base{32}{5}$ ;
$\base{31130}{4} = 860 = \base{1011212}{3}$ ;
$\base{134334}{5} = 5\,594 = \base{22211}{7}$ ;
$\base{203102}{6} = 16\,238 = \base{24242}{9}$ ;
$\base{35006}{7} = 8\,924 = \base{105152}{6}$ ;
$\base{7215}{8} = 3\,725 = \base{322031}{4}$.
(Pour le dernier, on peut aussi passer par la base $2$ : $\base{7215}{8} = \base{111\,010\,001\,101}{2}$,
puis on regroupe les chiffres par deux : $\base{11\,10\,10\,00\,11\,01}{2} = \base{322031}{4}$.)
\end{sol}
%
\begin{sol}[Exercice~\ref{exo:arith-5}]
\begin{enumerate}
\item $(n-1) + n + (n+1) = 3n$ est divisible par $3$.
\item Parmi $n-1$, $n$, $n+1$, l'un est congru \`a $0$ modulo $3$ (les trois restes modulo $3$ sont $0$, $1$ et $2$ dans un certain ordre) ;
      le produit est donc divisible par $3$.
\end{enumerate}
\end{sol}
%
\begin{sol}[Exercice~\ref{exo:arith-6}]
Soit $p \geq 5$ premier et $r$ son reste modulo $6$. Si $r \in \{0, 2, 4\}$, $p$ est pair et $p > 2$ : impossible.
Si $r = 3$, $3$ divise $p$ et $p > 3$ : impossible. Donc $r = 1$ ou $r = 5$, c'est-\`a-dire
$p = 6k+1$ ou $p = 6k + 5 = 6(k+1) - 1$ ; comme $p \geq 5$, l'entier $k$ (ou $k+1$) est sup\'erieur ou \'egal \`a $1$.
\end{sol}
%
\begin{sol}[Exercice~\ref{exo:arith-7}]
\begin{enumerate}
\item $n = 16 q + r$ avec $r = q^2$ et $0 \leq q^2 < 16$, donc $q \in \{0, 1, 2, 3\}$ ; $q = 0$ donne $n = 0$, exclu.
      On obtient $n \in \{17, 36, 57\}$.
\item $n = 125 q + q^3$ avec $0 \leq q^3 < 125$, donc $q \in \{1, 2, 3, 4\}$ : $n \in \{126, 258, 402, 564\}$.
\end{enumerate}
\end{sol}
%
\begin{sol}[Exercice~\ref{exo:arith-8}]
\begin{enumerate}
\item On a $n^3 - 8 = (n-2)(n^2 + 2n + 4)$, donc $n^3 - 2 = (n-2)(n^2+2n+4) + 6$.
      Si $n - 2 > 6$, c'est-\`a-dire $n \geq 9$, c'est la division euclidienne et le reste est $6$.
      Si $3 \leq n \leq 8$, le reste est celui de $6$ par $n - 2$ :
      $0$ pour $n = 3, 4, 5, 8$ ; $2$ pour $n = 6$ ; $1$ pour $n = 7$.
\item $n - 2$ divise $n^3 - 2$ si et seulement si $n - 2$ divise $6$, soit $n - 2 \in \{1, 2, 3, 6\}$ : $n \in \{3, 4, 5, 8\}$.
\end{enumerate}
\end{sol}
%
\begin{sol}[Exercice~\ref{exo:arith-9}]
\begin{enumerate}
\item $2^{n+3} = 8 \cdot 2^n$ et $8 \equiv 1 \ [7]$, donc $2^{n+3} \equiv 2^n \equiv r \ [7]$ : le reste est encore $r$.
      Les restes se r\'ep\`etent donc de $3$ en $3$ ; avec $2^0 = 1$, $2^1 = 2$, $2^2 = 4$ :
      le reste de $2^n$ est $1$ si $n = 3k$, $2$ si $n = 3k+1$, $4$ si $n = 3k+2$.
\item $23 \equiv 2 \ [7]$ et $41 = 3 \times 13 + 2$, donc $23^{41} \equiv 2^{41} \equiv 4 \ [7]$. Le reste est $4$.
\end{enumerate}
\end{sol}
%
\begin{sol}[Exercice~\ref{exo:arith-10}]
\begin{enumerate}
\item Les restes de $2^1, \ldots, 2^6$ sont $2, 4, 1, 2, 4, 1$.
\item $2^{n+3} - 2^n = 2^n (8 - 1) = 7 \cdot 2^n$. Les restes se r\'ep\`etent donc avec la p\'eriode $3$, et $2^n \equiv 1 \ [7]$
      si et seulement si $n$ est un multiple de $3$.
\item $9 \equiv 2 \ [7]$, donc $9^n \equiv 2^n \ [7]$ : m\^eme reste.
\item $2^{999} \cdot 9^{222} \equiv 2^{999} \cdot 2^{222} = 2^{1221} \ [7]$ et $1221 = 3 \times 407$ ; le reste est $1$.
\end{enumerate}
\end{sol}
%
\begin{sol}[Exercice~\ref{exo:arith-11}]
\begin{enumerate}
\item $a$ prend $9$ valeurs ($1$ \`a $9$) et $b$ en prend $10$ : $E$ contient $90$ nombres.
\item $N = a\,10^5 + b\,10^4 + a\,10^3 + b\,10^2 + a\,10 + b = (10a + b)(10^4 + 10^2 + 1) = 10101 \times \overline{ab}$.
\end{enumerate}
\end{sol}
%
\begin{sol}[Exercice~\ref{exo:arith-12}]
\begin{enumerate}
\item $10^0 = 1$ : reste $1$. Pour $p \geq 1$ : $10 \equiv 4 \ [6]$ et $4^2 = 16 \equiv 4 \ [6]$, donc par r\'ecurrence $10^p \equiv 4 \ [6]$ : reste $4$.
\item $x = a_0 + \sum_{p=1}^n a_p 10^p \equiv a_0 + 4\,(a_1 + \cdots + a_n) \ [6]$.
      Donc $6 | x$ si et seulement si $6 | 4\,(a_n + \cdots + a_1) + a_0$.
\end{enumerate}
\end{sol}
%
\begin{sol}[Exercice~\ref{exo:arith-13}]
\begin{enumerate}
\item $N^2 = 4k^2 + 4k + 1 = 4k(k+1) + 1$. Comme $k(k+1)$ est pair (produit de deux entiers cons\'ecutifs), $8 | 4k(k+1)$ : le reste est $1$.
\item Si $x$ est pair, $x^2$ est pair alors que $8y + 1$ est impair : pas de solution. Si $x = 2k+1$, l'\'equation devient
      $4k(k+1) = 8y$, soit $y = \dfrac{k(k+1)}{2}$, qui est bien un entier.
      Les solutions sont les couples $\left( 2k+1 \;;\; \dfrac{k(k+1)}{2} \right)$, $k \in \mathbb{Z}$.
\item Un point $(x ; y)$ \`a coordonn\'ees enti\`eres est sur $(P)$ si et seulement si $x^2 = 8y + 1$ ;
      il y en a une infinit\'e (un par valeur de $k$), par exemple $(1;0)$, $(3;1)$, $(5;3)$, $(7;6)$.
\end{enumerate}
\end{sol}
%
\begin{sol}[Exercice~\ref{exo:arith-14}]
\begin{enumerate}
\item $12 \equiv 1 \ [11]$, donc $12^p \equiv 1 \ [11]$ et $N = \sum a_p 12^p \equiv \sum a_p \ [11]$.
\item Un nombre \'ecrit en base $12$ est divisible par $11$ si et seulement si la somme de ses chiffres l'est.
\item $N = 144x + 12y + 4 = 300 + 10y + x$, donc $143x + 2y = 296$. Comme $0 \leq 2y \leq 18$, on a $278 \leq 143x \leq 296$, donc $x = 2$ et $y = 5$.
      Alors $N = 352$ ; v\'erification : $\overline{254}^{_{(12)}} = 288 + 60 + 4 = 352$ et $2 + 5 + 4 = 11$ : $N$ est bien divisible par $11$ ($352 = 11 \times 32$).
\end{enumerate}
\end{sol}
%
\begin{sol}[Exercice~\ref{exo:arith-15}]
\begin{enumerate}
\item $2^6 - 1 = 63 = 7 \times 9$ ; $3^6 - 1 = 728 = 7 \times 104$ ; $4^6 - 1 = 4095 = 7 \times 585$ ; $5^6 - 1 = 15624 = 7 \times 2232$.
\item $A_{n+6} - A_n = 2^n(2^6 - 1) + 3^n(3^6-1) + 4^n(4^6-1) + 5^n(5^6-1)$ est une somme de multiples de $7$.
\item On a $A_{n+6} \equiv A_n \ [7]$ ; par r\'ecurrence sur $q$, $A_{6q + r} \equiv A_r \ [7]$.
\item $A_0 = 4$, $A_1 = 14$, $A_2 = 54$, $A_3 = 224$, $A_4 = 978$, $A_5 = 4424$, dont les restes modulo $7$ sont
      $4$, $0$, $5$, $0$, $5$, $0$. Donc $A_n$ est divisible par $7$ si et seulement si $r \in \{1, 3, 5\}$,
      c'est-\`a-dire si et seulement si $n$ est impair.
\end{enumerate}
\end{sol}
%
\begin{sol}[Exercice~\ref{exo:arith-16}]
Algorithme d'Euclide (suite des restes) :
\begin{enumerate}
\item $297\,234$, $202\,643$, $94\,591$, $13\,461$, $364$, $357$, $7$, $0$ : $\pgcd{499\,877}{797\,111} = 7$.
\item $10\,204$, $5\,935$, $4\,269$, $1\,666$, $937$, $729$, $208$, $105$, $103$, $2$, $1$, $0$ : PGCD $= 1$.
\item $6\,928$, $3\,169$, $590$, $219$, $152$, $67$, $18$, $13$, $5$, $3$, $2$, $1$, $0$ : PGCD $= 1$.
\item $2\,197$, $1\,973$, $224$, $181$, $43$, $9$, $7$, $2$, $1$, $0$ : PGCD $= 1$.
\item $4\,109$, $3\,127$, $982$, $181$, $77$, $27$, $23$, $4$, $3$, $1$, $0$ : PGCD $= 1$.
\end{enumerate}
\end{sol}
%
\begin{sol}[Exercice~\ref{exo:arith-17}]
Avec l'algorithme d'Euclide \'etendu :
\begin{enumerate}
\item $\pgcd{699}{624} = 3$ et $699 \times 25 + 624 \times (-28) = 17\,475 - 17\,472 = 3$.
\item $\pgcd{972}{962} = 2$ et $972 \times (-96) + 962 \times 97 = -93\,312 + 93\,314 = 2$.
\item $\pgcd{875}{788} = 1$ et $875 \times (-317) + 788 \times 352 = -277\,375 + 277\,376 = 1$.
\item $\pgcd{934}{414} = 2$ et $934 \times (-82) + 414 \times 185 = -76\,588 + 76\,590 = 2$.
\item $\pgcd{872}{432} = 8$ et $872 \times 1 + 432 \times (-2) = 8$.
\end{enumerate}
\end{sol}
%
\begin{sol}[Exercice~\ref{exo:arith-18}]
On applique la m\'ethode de la le\c{c}on (solution particuli\`ere, puis lemme de Gauss) ; $k$ d\'ecrit $\mathbb{Z}$.
\begin{enumerate}
\item $3 \times 2 + 5 \times (-1) = 1$ : $S = \{ (2 + 5k ; -1 - 3k) \}$.
\item $13 \times (-8) + 21 \times 5 = 1$ : $S = \{ (-8 + 21k ; 5 - 13k) \}$.
\item $37 \times (-17) + 45 \times 14 = 1$ : $S = \{ (-17 + 45k ; 14 - 37k) \}$.
\item $3 \times (-1) + 5 \times 1 = 2$ : $S = \{ (-1 + 5k ; 1 - 3k) \}$.
\item $13 \times (-3) + 21 \times 2 = 3$ : $S = \{ (-3 + 21k ; 2 - 13k) \}$.
\item $37 \times 22 + 45 \times (-18) = 4$ : $S = \{ (22 + 45k ; -18 - 37k) \}$.
\item $\pgcd{6}{10} = 2$ divise $6$ ; l'\'equation \'equivaut \`a $3u + 5v = 3$, et $3 \times 1 + 5 \times 0 = 3$ : $S = \{ (1 + 5k ; -3k) \}$.
\item $\pgcd{39}{63} = 3$ divise $6$ ; l'\'equation \'equivaut \`a $13u + 21v = 2$, et $13 \times 5 + 21 \times (-3) = 2$ :
      $S = \{ (5 + 21k ; -3 - 13k) \}$.
\end{enumerate}
Par exemple pour 1 : si $3u + 5v = 1 = 3 \times 2 + 5 \times (-1)$, alors $3(u - 2) = 5(-1 - v)$ ; $5$ divise $3(u-2)$ et
$\pgcd{5}{3} = 1$, donc $5 | u - 2$ (Gauss) : $u = 2 + 5k$, puis $v = -1 - 3k$ ; r\'eciproquement ces couples conviennent.
\end{sol}
%
\begin{sol}[Exercice~\ref{exo:arith-19}]
\begin{enumerate}
\item Soit $d$ un diviseur commun de $5a + 3b$ et $3a + 2b$. Alors $d$ divise $2(5a+3b) - 3(3a+2b) = a$
      et $5(3a+2b) - 3(5a+3b) = b$ ; donc $d$ divise $\pgcd{a}{b} = 1$. Les deux nombres sont premiers entre eux.
\item \begin{enumerate}
      \item $n^2 + 1 = (n+1)(n-1) + 2$ ; un diviseur commun de $n^2+1$ et $n+1$ divise donc $2$.
      \item Si $n$ est pair, $n+1$ est impair, donc le PGCD (qui divise $2$) vaut $1$.
            Si $n$ est impair, $n + 1$ et $n^2 + 1$ sont pairs : le PGCD vaut $2$ (par exemple $n = 3$ : $\pgcd{10}{4} = 2$).
      \end{enumerate}
\end{enumerate}
\end{sol}
%
\begin{sol}[Exercice~\ref{exo:arith-20}]
\begin{enumerate}
\item $123 = 7 \times 17 + 4$ et $456 = 7 \times 65 + 1$ : l'\'equation \'equivaut \`a $4x \equiv 1 \ [7]$.
      Comme $4 \times 2 = 8 \equiv 1 \ [7]$, on multiplie par $2$ : $x \equiv 2 \ [7]$. Les solutions sont les $x = 2 + 7k$, $k \in \mathbb{Z}$.
\item \begin{enumerate}
      \item $78 = 2 \times 3 \times 13$ et $14 = 2 \times 7$.
      \item $D_{78} = \{1, 2, 3, 6, 13, 26, 39, 78\}$ et $D_{14} = \{1, 2, 7, 14\}$.
      \item $g = 2$ et $m = 2 \times 3 \times 7 \times 13 = 546$ (on v\'erifie que $g \cdot m = 1092 = 78 \times 14$).
      \end{enumerate}
\end{enumerate}
\end{sol}
%
\begin{sol}[Exercice~\ref{exo:arith-21}]
\begin{enumerate}
\item Un nombre divisible par $14$ et par $34$ est un multiple de $PPCM(14;34) = 2 \times 7 \times 17 = 238$.
      Les multiples de $238$ ayant trois chiffres sont $238$, $476$, $714$ et $952$.
\item On \'ecrit $a = 6a'$, $b = 6b'$ avec $\pgcd{a'}{b'} = 1$ ; alors $36\, a'b' = 2700$, soit $a'b' = 75 = 3 \times 5^2$.
      Les couples $(a';b')$ premiers entre eux de produit $75$ sont $(1;75)$, $(3;25)$, $(25;3)$, $(75;1)$, d'o\`u
      $(a;b) \in \{ (6;450), (18;150), (150;18), (450;6) \}$.
      Dans tous les cas, $PPCM(a;b) = \dfrac{ab}{\pgcd{a}{b}} = \dfrac{2700}{6} = 450$.
\end{enumerate}
\end{sol}
%
\begin{sol}[Exercice~\ref{exo:arith-22}]
\begin{enumerate}
\item Pour $n = 0$ : $9^0 - 2^0 = 0$, divisible par $7$. Si $7 | 9^k - 2^k$, alors
      $9^{k+1} - 2^{k+1} = 9\,(9^k - 2^k) + 7 \cdot 2^k$ est divisible par $7$.
\item $5(3n - 2) - 3(5n - 3) = -1$ ; un diviseur commun divise donc $1$ : les deux nombres sont premiers entre eux
      (c'est aussi une cons\'equence du th\'eor\`eme de B\'ezout).
\item \begin{enumerate}
      \item $\bar{5}^2 - \bar{5} + \bar{1} = \bar{21} = \bar{0}$.
      \item $(x - \bar{4})^2 - \bar{1} = x^2 - \bar{8}x + \bar{15} = x^2 - x + \bar{1}$ (car $\bar{8} = \bar{1}$ et $\bar{15} = \bar{1}$) ;
            donc $(E) \iff (x - \bar{4})^2 = \bar{1}$.
      \item $(x-\bar{4})^2 = \bar{1} \iff (x - \bar{4} - \bar{1})(x - \bar{4} + \bar{1}) = \bar{0} \iff (x - \bar{5})(x - \bar{3}) = \bar{0}$.
            Comme $7$ est premier, un produit est nul dans $\mathbb{Z}/7\mathbb{Z}$ si et seulement si l'un des facteurs l'est :
            $S = \{ \bar{3} ; \bar{5} \}$.
      \end{enumerate}
\end{enumerate}
\end{sol}
%
\begin{sol}[Exercice~\ref{exo:arith-23}]
\begin{enumerate}
\item $319 = 11 \times 29$, $232 = 8 \times 29$, $145 = 5 \times 29$ : $PGCD(319, 232, 145) = 29$.
      L'\'equation \'equivaut \`a $11x + 8y = 5$. Solution particuli\`ere : $11 \times (-1) + 8 \times 2 = 5$.
      Alors $11(x + 1) = 8(2 - y)$, $8$ divise $11(x+1)$ et $\pgcd{8}{11} = 1$, donc $x = -1 + 8k$ et $y = 2 - 11k$, $k \in \mathbb{Z}$.
\item Solution particuli\`ere : $13 \times 3 - 7 \times 4 = 11$. Alors $13(x - 3) = 7(y - 4)$, donc $x = 3 + 7k$, $y = 4 + 13k$, $k \in \mathbb{Z}$.
      La condition $x \geq 0$ et $y \geq 0$ impose $k \geq 0$ (pour $k = -1$, $x = -4$).
      $S = \{ (3 + 7k ; 4 + 13k) \mid k \in \mathbb{N} \}$.
\end{enumerate}
\end{sol}
%
\begin{sol}[Exercice~\ref{exo:arith-24}]
Modulo $3$ : si $n \equiv 0$, $n^2 + n + 2 \equiv 2$ ; si $n \equiv 1$, $n^2 + n + 2 \equiv 4 \equiv 1$ ; si $n \equiv 2$, $n^2 + n + 2 \equiv 8 \equiv 2$.
Le reste n'est jamais $0$.
\end{sol}
%
\begin{sol}[Exercice~\ref{exo:arith-25}]
\begin{enumerate}
\item \begin{enumerate}
      \item $N = 2 \cdot 81 + 1 \cdot 27 + 0 \cdot 9 + 1 \cdot 3 + 2 = 194$.
      \item $100 = 3 \times 33 + 1$, $33 = 3 \times 11 + 0$, $11 = 3 \times 3 + 2$, $3 = 3 \times 1 + 0$, $1 = 3 \times 0 + 1$ : $100 = \base{10201}{3}$.
      \item $N = a_0 + 3\,(a_1 + 3a_2 + \cdots)$, donc $N \equiv a_0 \ [3]$ : $N$ est divisible par $3$ si et seulement si son dernier chiffre est $0$.
            De m\^eme $N \equiv a_0 + 3a_1 \ [9]$ et $0 \leq a_0 + 3a_1 \leq 8$ : $N$ est divisible par $9$ si et seulement si ses deux derniers chiffres sont $0$.
      \end{enumerate}
\item \begin{enumerate}
      \item 
      $\begin{array}{|c|c|c|c|}
      \hline + & \bar{0} & \bar{1} & \bar{2} \\ \hline
      \bar{0} & \bar{0} & \bar{1} & \bar{2} \\ \hline
      \bar{1} & \bar{1} & \bar{2} & \bar{0} \\ \hline
      \bar{2} & \bar{2} & \bar{0} & \bar{1} \\ \hline
      \end{array}
      \qquad
      \begin{array}{|c|c|c|c|}
      \hline \times & \bar{0} & \bar{1} & \bar{2} \\ \hline
      \bar{0} & \bar{0} & \bar{0} & \bar{0} \\ \hline
      \bar{1} & \bar{0} & \bar{1} & \bar{2} \\ \hline
      \bar{2} & \bar{0} & \bar{2} & \bar{1} \\ \hline
      \end{array}$
      \item Pour $x = \bar{0}, \bar{1}, \bar{2}$, $x^2 + x + \bar{2}$ vaut respectivement $\bar{2}$, $\bar{4} = \bar{1}$, $\bar{8} = \bar{2}$ : jamais $\bar{0}$.
      \item $x^2 + \bar{2}x + \bar{1} = (x + \bar{1})^2$ ; il est nul si et seulement si $x + \bar{1} = \bar{0}$ ($3$ est premier), soit $x = \bar{2}$.
            (V\'erification : $\bar{4} + \bar{4} + \bar{1} = \bar{9} = \bar{0}$.)
      \item En soustrayant la premi\`ere \'equation de la seconde : $y = \bar{0} - \bar{2} = \bar{1}$, puis $x = \bar{2} - \bar{1} = \bar{1}$.
            Solution unique : $(\bar{1} ; \bar{1})$.
      \end{enumerate}
\end{enumerate}
\end{sol}
%
\begin{sol}[Exercice~\ref{exo:arith-26}]
\begin{enumerate}
\item Pour $x = \bar{0}, \bar{1}, \bar{2}, \bar{3}, \bar{4}$ : $x^2 = \bar{0}, \bar{1}, \bar{4}, \bar{4}, \bar{1}$.
\item $x^2 = \bar{4} \iff x \in \{ \bar{2} ; \bar{3} \}$ ; $x^2 = \bar{2}$ n'a pas de solution.
      Si $n^2$ se terminait par $2$ ou $7$, on aurait $n^2 \equiv 2 \ [5]$ ; s'il se terminait par $3$ ou $8$, on aurait $n^2 \equiv 3 \ [5]$.
      Or un carr\'e est congru \`a $0$, $1$ ou $4$ modulo $5$ : c'est impossible.
\item $x^4 = (x^2)^2$ vaut $\bar{0}$ pour $x = \bar{0}$ et $\bar{1}$ sinon. Donc $x^5 = x \cdot x^4 = x$ pour $x \neq \bar{0}$, et $\bar{0}^5 = \bar{0}$.
\item Pour tout entier $n$, $\overline{n}^5 = \overline{n}$ dans $\mathbb{Z}/5\mathbb{Z}$, c'est-\`a-dire $n^5 \equiv n \ [5]$ : $5 | n^5 - n$.
\item \begin{enumerate}
      \item $n^5 - n = n(n^4 - 1) = n(n^2 - 1)(n^2 + 1) = (n-1)\,n\,(n+1)\,(n^2+1)$.
      \item Parmi les trois entiers cons\'ecutifs $n-1$, $n$, $n+1$, l'un est pair et l'un est multiple de $3$ :
            $2$ et $3$ divisent $n^5 - n$, et comme $\pgcd{2}{3} = 1$, $6$ divise $n^5 - n$.
            Enfin $5$ et $6$ divisent $n^5 - n$ et $\pgcd{5}{6} = 1$, donc $30$ divise $n^5 - n$.
      \end{enumerate}
\item $10$ divise $n^5 - n$, donc $n^5 \equiv n \ [10]$ : $n^5$ et $n$ ont le m\^eme chiffre des unit\'es
      (par exemple $7^5 = 16\,807$).
\end{enumerate}
\end{sol}
%
\begin{sol}[Exercice~\ref{exo:arith-27}]
Oui, car $\pgcd{5}{3} = 1$ et $2 \times 3 - 1 \times 5 = 1$. On remplit le r\'ecipient de $3$ litres et on le verse dans celui de $5$ litres ;
on remplit de nouveau le r\'ecipient de $3$ litres et on le verse dans celui de $5$ litres jusqu'\`a ce qu'il soit plein :
on n'y verse que $2$ litres, et il reste exactement $1$ litre dans le r\'ecipient de $3$ litres.
\end{sol}
%
\begin{sol}[Exercice~\ref{exo:arith-28}]
\begin{enumerate}
\item Il faut trouver des entiers naturels $x$ et $y$ avec $4x - 10y = 18$, soit $2x - 5y = 9$ :
      $x = 7$, $y = 1$ convient. Andry donne $7$ pi\`eces de $4$~Ar ($28$~Ar) et Bema lui rend une pi\`ece de $10$~Ar : c'est possible.
\item Il faudrait $10x - 4y = 25$ ; le membre de gauche est pair et $25$ est impair : c'est impossible
      ($\pgcd{10}{4} = 2$ ne divise pas $25$).
\end{enumerate}
\end{sol}
%
\begin{sol}[Exercice~\ref{exo:arith-29}]
$10^{10}$ se termine par $0$. Pour $k \geq 1$, $5^k$ se termine par $5$ (car $5 \times 5 = 25$), donc $5^{5^5}$ se termine par $5$.
Pour $k \geq 1$, $2^{k+4} - 2^k = 15 \cdot 2^k$ est divisible par $10$ : le dernier chiffre de $2^k$ ne d\'epend que du reste de $k$ modulo $4$
($2$, $4$, $8$, $6$ pour $k \equiv 1, 2, 3, 0$). Or $3^2 = 9 \equiv 1 \ [4]$, donc $3^{4^5} = (3^2)^{512} \equiv 1 \ [4]$ :
$2^{3^{4^5}}$ se termine par $2$.
\end{sol}
%
\begin{sol}[Exercice~\ref{exo:arith-30}]
$\pgcd{4}{12} = 4$ ; $\pgcd{60}{24} = 12$ ;
$9379 = 1 \cdot 8039 + 1340$, $8039 = 5 \cdot 1340 + 1339$, $1340 = 1 \cdot 1339 + 1$ : $\pgcd{9379}{8039} = 1$.
\end{sol}
%
\begin{sol}[Exercice~\ref{exo:arith-31}]
On travaille modulo $7$ : revenir $N$ jours apr\`es un dimanche, c'est tomber le jour << dimanche $+ (N \bmod 7)$ >>.
\begin{enumerate}
\item $40 = 7 \times 5 + 5$ : vendredi. $40\,000 \equiv 2 \ [7]$ : mardi. $40\,000\,000 \equiv 5 \ [7]$ : vendredi.
\item On remonte le temps : $50 \equiv 1 \ [7]$, le voyage a commenc\'e un samedi.
      $50\,000 \equiv 6 \ [7]$ : dimanche $- 6$ jours, soit un lundi. $50\,000\,000 \equiv 1 \ [7]$ : un samedi.
\end{enumerate}
\end{sol}
%
\begin{sol}[Exercice~\ref{exo:arith-32}]
Chaque soir, le propri\'etaire d\'etient un ensemble de morceaux dont la longueur totale est le num\'ero du jour.
Avec $k$ morceaux, il n'y a que $2^k - 1$ ensembles non vides de morceaux ; il en faut au moins $31$ diff\'erents, donc $2^k - 1 \geq 31$ et $k \geq 5$.
Cinq morceaux suffisent : $1$, $2$, $4$, $8$ et $16$~cm (total $31$~cm). Le jour $n$, le propri\'etaire d\'etient les morceaux
correspondant aux chiffres $1$ de l'\'ecriture de $n$ en base $2$ (par exemple $13 = \base{1101}{2}$ : morceaux de $8$, $4$ et $1$~cm).
\end{sol}
%
\begin{sol}[Exercice~\ref{exo:arith-33}]
\begin{enumerate}
\item Si $d$ divise $n+7$ et $2n+5$, il divise $2(n+7) - (2n+5) = 9$.
\item $n + 7 \equiv n - 2 \ [3]$, donc $3 | n+7 \iff n \equiv 2 \ [3]$. Dans ce cas $2n + 5 \equiv 4 + 5 = 9 \equiv 0 \ [3]$.
\item Le PGCD de $n+7$ et $2n+5$ divise $9$ : il vaut $1$, $3$ ou $9$. Il est diff\'erent de $1$ si et seulement si $3$ divise les deux nombres,
      c'est-\`a-dire si et seulement si $n \equiv 2 \ [3]$. La fraction est irr\'eductible si et seulement si $n \not\equiv 2 \ [3]$.
\item $9 | n + 7 \iff n \equiv 2 \ [9]$, et alors $2n + 5 \equiv 9 \equiv 0 \ [9]$. On simplifie par $9$ exactement lorsque $n \equiv 2 \ [9]$.
      Pour $n = 20$ ($20 \equiv 2 \ [9]$) : $\dfrac{27}{45} = \dfrac{3}{5}$.
\end{enumerate}
\end{sol}
%
\begin{sol}[Exercice~\ref{exo:arith-34}]
\begin{enumerate}
\item Le nombre obtenu est $10x + 1$. La condition $10x + 1 = 3(10^k + x)$ s'\'ecrit $7x = 3 \cdot 10^k - 1$.
\item $10 \equiv 3 \ [7]$ ; les restes de $10^1, \ldots, 10^6$ sont $3$, $2$, $6$, $4$, $5$, $1$.
\item Comme $3 \times 5 = 15 \equiv 1 \ [7]$, on a $3 \cdot 10^k \equiv 1 \ [7] \iff 10^k \equiv 5 \ [7]$ (on multiplie par $5$, ou par $3$ pour revenir).
\item Le plus petit $k$ est $k = 5$ ; alors $x = \dfrac{299\,999}{7} = 42\,857 < 10^5$ et $\boxed{N = 142\,857}$.
      V\'erification : $3 \times 142\,857 = 428\,571$.
\end{enumerate}
\end{sol}
%
\begin{sol}[Exercice~\ref{exo:arith-35}]
\begin{enumerate}
\item $5^0 = 1$, $5^1 = 5$, $5^2 = 25 \equiv 12$, $5^3 \equiv 60 \equiv 8$, $5^4 \equiv 40 \equiv 1 \ [13]$.
      Comme $5^{n+4} \equiv 5^n \ [13]$, le reste de $5^n$ est $1$, $5$, $12$, $8$ selon que $n \equiv 0, 1, 2, 3 \ [4]$.
\item $13 | 5^n - 1 \iff n \equiv 0 \ [4]$ ; $13 | 5^n + 1 \iff 5^n \equiv 12 \ [13] \iff n \equiv 2 \ [4]$.
\item Il faudrait $5^n \equiv -2 \equiv 11 \ [13]$, or $11$ n'est pas dans la liste $\{1, 5, 12, 8\}$.
\item $2026 \equiv 2$ et $2027 \equiv 3 \ [4]$, donc $5^{2026} + 5^{2027} \equiv 12 + 8 = 20 \equiv 7 \ [13]$ : le reste est $7$.
\end{enumerate}
\end{sol}
%
\begin{sol}[Exercice~\ref{exo:arith-36}]
\begin{enumerate}
\item Les \'el\'ements de $A$ sont les $3 + 5p$ et ceux de $B$ les $7 + 8q$, avec $0 \leq p, q \leq 99$.
\item $5 \times 4 - 8 \times 2 = 4$. Si $5p - 8q = 4$, alors $5(p - 4) = 8(q - 2)$, donc $8 | p - 4$ (Gauss) :
      $p = 4 + 8t$, $q = 2 + 5t$, $t \in \mathbb{Z}$ ; r\'eciproquement ces couples conviennent.
\item Les nombres communs sont les $3 + 5(4 + 8t) = 23 + 40t$ : suite arithm\'etique de premier terme $23$ et de raison $40 = PPCM(5;8)$.
\item Il faut $0 \leq 4 + 8t \leq 99$ et $0 \leq 2 + 5t \leq 99$, soit $0 \leq t \leq 11$ : $A \cap B$ contient $12$ nombres
      ($23$, $63$, \ldots, $463$), et $A \cup B$ contient $100 + 100 - 12 = 188$ nombres distincts.
\end{enumerate}
\end{sol}
%
\begin{sol}[Exercice~\ref{exo:arith-37}]
\begin{enumerate}
\item $a | b$ et $a | c$ (car $a | b$ et $b | c$), donc $a | a + b + c = 55$.
\item $b > a$ donne $k \geq 2$ ; $c > b$ donne $m \geq 2$. Et $a + ka + kma = a(1 + k + km) = 55$.
\item $a \in \{1, 5, 11, 55\}$.
      \begin{itemize}
      \item $a = 1$ : $k(1 + m) = 54$ avec $k \geq 2$ et $1 + m \geq 3$ : $k \in \{2, 3, 6, 9, 18\}$, d'o\`u
            $(1, 2, 52)$, $(1, 3, 51)$, $(1, 6, 48)$, $(1, 9, 45)$, $(1, 18, 36)$.
      \item $a = 5$ : $k(1+m) = 10$, donc $k = 2$, $m = 4$ : $(5, 10, 40)$.
      \item $a = 11$ : $k(1+m) = 4$ est impossible avec $k \geq 2$ et $m \geq 2$ ; $a = 55$ est impossible.
      \end{itemize}
      Il y a six triplets.
\item La plus grande valeur de $c$ est $52$.
\end{enumerate}
\end{sol}
%
\begin{sol}[Exercice~\ref{exo:arith-38}]
\begin{enumerate}
\item $(n+3)(n-3) + 16 = n^2 - 9 + 16 = n^2 + 7$.
\item $n + 3$ divise $(n+3)(n-3)$ ; il divise donc $n^2 + 7$ si et seulement s'il divise $n^2 + 7 - (n+3)(n-3) = 16$.
\item $n + 3 \in \{\pm 1, \pm 2, \pm 4, \pm 8, \pm 16\}$, soit $n \in \{-19, -11, -7, -5, -4, -2, -1, 1, 5, 13\}$.
\item Le plus grand est $n = 13$ : $13^2 + 7 = 176 = 16 \times 11$, le quotient est $11$.
\end{enumerate}
\end{sol}
%
\begin{sol}[Exercice~\ref{exo:arith-39}]
\begin{enumerate}
\item $\dfrac{1}{12}$, $\dfrac{5}{12}$, $\dfrac{7}{12}$, $\dfrac{11}{12}$ ; leur somme est $\dfrac{24}{12} = 2$.
\item Un entier divise $k$ et $n$ si et seulement s'il divise $n - k$ et $n$ ; les deux couples ont les m\^emes diviseurs communs, donc le m\^eme PGCD.
\item $\dfrac{k}{n}$ irr\'eductible $\iff \pgcd{k}{n} = 1 \iff \pgcd{n-k}{n} = 1 \iff \dfrac{n-k}{n}$ irr\'eductible.
      Si les deux fractions \'etaient \'egales, on aurait $n = 2k$, puis $\pgcd{k}{n} = k = 1$ et $n = 2$, ce qui est exclu.
\item Les fractions irr\'eductibles se r\'epartissent en paires $\left\lbrace \dfrac{k}{n} ; \dfrac{n-k}{n} \right\rbrace$ de somme $1$.
      Il y a donc $\dfrac{N}{2}$ paires (donc $N$ est pair) et $S = \dfrac{N}{2} \times 1 = \dfrac{N}{2}$.
\item Pour $n = 15 = 3 \times 5$, les $k$ premiers avec $15$ sont $1, 2, 4, 7, 8, 11, 13, 14$ : $N = 8$ et $S = 4$.
\end{enumerate}
\end{sol}
%
\begin{sol}[Exercice~\ref{exo:arith-40}]
\begin{enumerate}
\item Posons $x = 2^q$. On a $x^p - 1 = (x - 1)(x^{p-1} + x^{p-2} + \cdots + x + 1)$, donc $2^q - 1$ divise $2^{pq} - 1$.
\item Supposons $n$ non premier. Si $n = 1$, $2^1 - 1 = 1$ n'est pas premier. Sinon $n = pq$ avec $1 < q < n$ ;
      alors $2^q - 1$ divise $2^n - 1$ et $1 < 2^q - 1 < 2^n - 1$ : $2^n - 1$ n'est pas premier.
      Par contrapos\'ee, si $2^n - 1$ est premier, $n$ est premier.
\item $p = 11$ : $2^{11} - 1 = 2047 = 23 \times 89$.
\end{enumerate}
\end{sol}
%
\begin{sol}[Exercice~\ref{exo:arith-42}]
\begin{enumerate}
\item \begin{enumerate}
      \item Par r\'ecurrence : $F_0 = 3 = F_1 - 2$ ($F_1 = 5$). Si $F_0 \cdots F_{n-1} = F_n - 2 = 2^{2^n} - 1$, alors
            $F_0 \cdots F_n = (2^{2^n} - 1)(2^{2^n} + 1) = 2^{2^{n+1}} - 1 = F_{n+1} - 2$.
      \item Soit $m < n$ et $d$ un diviseur commun de $F_m$ et $F_n$. Comme $F_m$ est un facteur de $F_0 \cdots F_{n-1} = F_n - 2$,
            $d$ divise $F_n - (F_n - 2) = 2$. Les nombres de Fermat sont impairs, donc $d = 1$.
      \end{enumerate}
\item $2^{2^n} \equiv -1 \ [F_n]$, donc $2^{2^{n+1}} \equiv 1 \ [F_n]$.
      Or $F_n - 1 = 2^{2^n}$ est un multiple de $2^{n+1}$ (car $2^n \geq n + 1$), donc $2^{F_n - 1} \equiv 1 \ [F_n]$
      et $2^{F_n} \equiv 2 \ [F_n]$ : $F_n$ divise $2^{F_n} - 2$.
\end{enumerate}
\end{sol}
%
\begin{sol}[Exercice~\ref{exo:arith-43}]
\begin{enumerate}
\item \begin{enumerate}
      \item $a = 2k$ et $3^a = 9^k \equiv 1^k = 1 \ [4]$.
      \item $2^{4k} = 16^k \equiv 1 \ [5]$.
      \item D'apr\`es (a), $3^a = 4k + 1$ avec $k \in \mathbb{N}$, donc $2^{3^a} = 2 \cdot 2^{4k} \equiv 2 \ [5]$ d'apr\`es (b).
            Ainsi $2^{3^a} - 2$ est divisible par $5$, et il est pair : il est divisible par $10$.
      \item $4^5$ est pair, donc $2^{3^{4^5}} \equiv 2 \ [10]$ : le dernier chiffre est $2$.
      \end{enumerate}
\item $7^4 = 2401 \equiv 1 \ [100]$. L'exposant $E = 7^{7^{\cdot^{\cdot^{7}}}}$ est une puissance de $7$ d'exposant impair, et $7 \equiv -1 \ [4]$,
      donc $E \equiv -1 \equiv 3 \ [4]$. Ainsi $7^E \equiv 7^3 = 343 \equiv 43 \ [100]$ : les deux derniers chiffres sont $43$.
\item Le nombre est divisible par $2^3 \times 5^3 = 1000$ (les deux exposants d\'epassent $3$) : ses trois derniers chiffres sont $000$.
\end{enumerate}
\end{sol}
%
\begin{sol}[Exercice~\ref{exo:arith-44}]
Notons $T_k$ le nombre form\'e de $k$ chiffres $3$. Parmi les $2020$ nombres $T_1, \ldots, T_{2020}$, deux ont le m\^eme reste
dans la division par $2019$ (il n'y a que $2019$ restes possibles) : $T_m \equiv T_k \ [2019]$ avec $k < m$.
Alors $2019$ divise $T_m - T_k = T_{m-k} \times 10^k$. Or $2019 = 3 \times 673$ est premier avec $10^k$, donc (Gauss)
$2019$ divise $T_{m-k}$. (Le plus petit tel terme a $224$ chiffres.)
\end{sol}
%
\begin{sol}[Exercice~\ref{exo:arith-45}]
$2019$ est impair, donc l'un des deux nombres est pair, c'est-\`a-dire \'egal \`a $2$ : $x = 2$, $y = 2017$.
$2017$ est premier (aucun nombre premier inf\'erieur ou \'egal \`a $43$ ne le divise et $\sqrt{2017} < 45$). Le seul couple est $(2 ; 2017)$.
\end{sol}
%
\begin{sol}[Exercice~\ref{exo:arith-46}]
\begin{enumerate}
\item Par exemple $P(40) = 1681 = 41^2$, $P(41) = 41 \times 43$, $P(44) = 2021 = 43 \times 47$, $P(49) = 2491 = 47 \times 53$,
      $P(81) = 6683 = 41 \times 163$. (Plus g\'en\'eralement, $41$ divise $P(41k)$.)
\item $p^2 - 1 = (p-1)(p+1)$. Comme $p$ est impair, $p - 1$ et $p + 1$ sont deux entiers pairs cons\'ecutifs, dont l'un est multiple de $4$ :
      $8$ divise $p^2 - 1$. Parmi $p - 1$, $p$, $p + 1$, l'un est multiple de $3$, et ce n'est pas $p$ : $3$ divise $p^2 - 1$.
      Comme $\pgcd{3}{8} = 1$, $24$ divise $p^2 - 1$.
\end{enumerate}
\end{sol}
\end{corrige}

